\documentclass[a4paper,11pt]{article}
\usepackage{longtable}
% 数式
\usepackage{amsmath,amsfonts,hyperref,amssymb,empheq}
\usepackage{latexsym,mathtools,array}
\usepackage{physics}
\usepackage{bm}
% 画像
\usepackage[dvipdfmx]{graphicx}
\usepackage{listings}
\usepackage{plistings}
\usepackage{booktabs}
\usepackage{tabularx}
\usepackage{tikz-cd}

% font
\usepackage[T1]{fontenc}
\usepackage[utf8]{inputenc}   % pdfLaTeX assumed; XeLaTeX/LuaLaTeX would drop the unicode
\usepackage{amssymb}
\usepackage{upgreek}
\usepackage{color}

 % Colors referenced by lstlean.tex -- all must be defined or compilation fails
\definecolor{keywordcolor}{rgb}{0.7, 0.1, 0.1}
\definecolor{tacticcolor}{rgb}{0.0, 0.1, 0.6}
\definecolor{commentcolor}{rgb}{0.4, 0.4, 0.4}
\definecolor{symbolcolor}{rgb}{0.0, 0.1, 0.6}
\definecolor{sortcolor}{rgb}{0.1, 0.5, 0.1}
\definecolor{attributecolor}{rgb}{0.7, 0.1, 0.1}
 
\def\lstlanguagefiles{lstlean.tex}
\lstset{language=lean}
\lstset{basicstyle=\ttfamily\small, breaklines=true,
        columns=fullflexible, showstringspaces=false}

\makeatletter % some generic helpers
\newcommand{\LCASES}[1]{$\m@th\displaystyle{#1}$\hfil}
\newcommand{\CCASES}[1]{\hfil$\m@th\displaystyle{#1}$\hfil}
\newcommand{\RCASES}[1]{\hfil$\m@th\displaystyle{#1}$}
\makeatother
\newcases{ecases}{\quad}{\CCASES{##}}{\LCASES{##}}{\lbrace}{.}
\newcases{ecases*}{\quad}{\CCASES{##}}{{##}\hfil}{\lbrace}{.}

\begin{document}
\title{Transformers as Functional Dynamics, equivalency between lambda calculas, linear logic, and composition of categories}
\author{Yoshinori Watanabe \\ \href{mailto::xiangzey@gmail.com}{xiangze@gmail.com}}
\date{\today}
\maketitle

\begin{abstract}
We give a compositional categorical account of a single Transformer layer. 
The value path of self-attention is a morphism of the symmetric monoidal closed category (SMCC) \((\mathbf{Vect},\otimes,\multimap)\),
whose internal language is multiplicative intuitionistic linear logic (MILL); under a linear approximation this lets us read in-context
computation as a linear \(\lambda\)-term. 
The MLP acts as a non-linear realization map and, empirically, as a key--value memory. The residual connection supplies an additive copy (the biproduct diagonal), which is a legitimate but weaker notion of copy than the forbidden multiplicative diagonal and than the exponential modality. The softmax attention matrix is a morphism of a Markov category --- the Kleisli category of the distribution monad --- and value mixing is the action of the associated \(D\)-algebra (expectation). 
Consequently a Transformer layer is naturally written as a composite of a Markov category and an SMCC, plus an additive/non-linear layer, rather than as any single closed category.
\end{abstract} 

\hypertarget{introduction}{%
\section{Introduction}\label{introduction}}

Almost universal computation power of Transfomers attracta many reseachers. Especiallty universal turing machine(UTM) ,lambda calculation theory and category theory are usually use to explain them.

Cartesian closed category(CCC) is defined having objects called direct product between two objects \(X\times Y\) and exponential object(ofnen
written \(Z^Y\)) .Intuitively an exponential object\(X^Y\) is set of all
morphisms from X to Y. has natural transformation \(Hom(X\times Y,X)\simeq Hom(X,Z^Y)\) for objects \(X,Y,Z\).

In \(\lambda\) calculas or programming points of view, morphism \(X \rightarrow Z^Y\) is currying, \(Z^Y \rightarrow Z\) is \(\lambda\) calculation, i.e eval of S-expression. As explained bellow eval is the oparation to generate attention matricx from matrix Q and K in transformers, function
application \(f\cdot f\) in functional dynamics.

But category \(\bf{Vect}\) whose objects is vector space, morphisms are linear transformation can not have nonlinear diagnal product. This is
not CCC but called Symmetric Monoid Closed Category(SMCC). In this categorey usual eval can not used for functions without constraint , but
use linear logic deduction, which treat propositions as finite resources when it used comsumed.

In same motivation research as this paper\cite{Topos},The assume Neural networks on category \(\bf{RELU}\) whose objects are usual vector space,
morphisms are partially linear map, because Relu is usually useed as activation function. Then transformers can be treated topos, which is
special case of CCC,and have univarsal higher order caluculation ability.

Linear logic is related to programming language such as Rust\cite{Rust} which constraint resource(such as variables) usage at one time. This reduces programming bugs. Also There is a research to connect linear logic and probablistic programming, bayesian inference\cite{Murfet2025}.

In this paper almost explaining about the relation between Transformers,
lambda calculation and category theory. But the original idea and motivation about higher order fuctions of Vector space and moprphism or
category of metafunction(funcsions between funcions) is from dynamical system of functions, this is deeply related to learnability of transformers. Some theorems explained following chapter are written and proven in Lean.

\hypertarget{contributions-of-this-paper}{%
\subsubsection{Contributions of this paper}\label{contributions-of-this-paper}}

\begin{itemize}
\item
  Points out the equivalence of self-attention and functional dynamics,
  property category .
\item
  Explains the relation between self-attention and symmetric monoidal
  closed category (SMCC) \((\mathbf{Vect},\otimes,\multimap)\)
  eval-apply loop which is required for in-context learning and the
  correspondence between markov category and attention matrix.
\item
  Numerical experiment about MLP function is only Key-value retirival or
  not. Comparison with different architectures (RNN,SSM) and ablation to
  find which components are responsible for this performance.
\item
  The proofs of formalizations are witten and proven in lean
\end{itemize}

Reading the residual stream as a linear polynomial extension \(\mathcal{K}[x]\), the ``apply'' of the loop is derived, not posited:
it is the evaluation counit of the abstraction--application adjunction that the (linear) deduction theorem produces from the polynomial
extension in the sense of Lambek and Došen. We ground each component in the existing categorical foundations of machine learning, embed results
from a controlled experiment that localizes the copy structure in softmax and the applied value in the MLP, and leave the Function-Vector
experiments as clearly marked placeholders to be filled from a real language model. On expressivity we defer to the topos analysis of
Villani and McBurney\cite{Topos}, whose piecewise-linear base is for ReLU networks and whose choose-eval decomposition independently reaches the
eval--apply reading; we measure a different quantity, resource structure, which a cartesian model cannot see. In the linear/PL-input
regime the value path --- residual copy included --- is expressible in \(!\)-free intuitionistic linear \(\lambda\)-calculus, and the hypothesis that the model performs only such resource-linear computation cannot be excluded; our experiments support rather than merely fail to refute it.

\hypertarget{preriminalies}{
\section{Preriminalies}\label{preriminalies}}

\hypertarget{attention-mechanism-transformers-and-their-components}{%
\subsubsection{Attention mechanism, Transformers and their Components}\label{attention-mechanism-transformers-and-their-components}}

Transformers are consists of several components, attention, MLP(multi layer perceptron ,FFN), softmax, residual connecctions and layer normalizations\cite{Attention}.

\begin{figure}
\centering
\includegraphics[width=1.2\textwidth]{img/Transformer.pdf}
\caption{Schematic architecture of a Transfomer}
\label{fig:transformer}
\end{figure}

Attentions are product of and input vector x. MLP is composition of all-to-all vector product using matrix product and activation function 
such as Relu or softmax. softmax function is usually used tu make attention matrix in conrast of Relu in the tail of MLP. Residual
connection (Resnet)\cite{Resnet} is often used in LLM. The benefits of Resnet are not only preserving information of earlier layers during training and
inference, but simplify loss landscape. Resnet with nearly identical matrix convertion are similar to differential operations which is called
neural ordinal differencial equiation(neural ODE). Layer normalizations are another important part of transformers to regulize internal data.

\begin{figure}
\centering
\includegraphics[width=.9\columnwidth]{img/summary.pdf}
\caption{architecture of attention mechanism}
\label{fig:summary}
\end{figure}

Positional encoders are also important for identify the order of tokens
which is encoded and put in attention mechanism. Layered transformers is
usually called large language models(LLM). LLMs have in-context learning
ability\cite{ICL}\cite{ICLgd} and scalability of learning. LLMs and their
variants made various applications and theoretical explanations.

Function vector(FV) \cite{FV} a concept embedded in LLM as a head of
transformer. FV is portable among but layer of LLM.

\hypertarget{functional-dynamics}{%
\subsubsection{Functional Dynamics}\label{functional-dynamics}}

Functional dynamics (FD)\cite{FD1} is introduced by function of
1-dimentional function (metafunction). The original form FD is define as following

\(f_{t+1}(x)=(f_t\cdot f_t)(x) +\epsilon f_t(x)\)

or

\(f_{t+1}(x)=(f_t\cdot g_t)(x) +\epsilon f_t(x)\)

Generally, the dynamics of functions are governed by fixed points and
hieralchy of fixed points and the structure complex behavior depends on
initial function f and parameter \(\epsilon\)\cite{FD2}. Regarding a
function as a graph drawn in 2D rectangle, function applicaton to other
function (\(f\cdot g\)) is described as matrix multiplication. In case attention mechanism, f and g corresponding to the non-zero elements of the matrix,  row is x -axis, column is y-axis graph.Then a matrix not only represent 1-dimentional function graph but

As a example \(d=d_q=d_k=d_v\) for simplicity, x has only 1 nonzero
value per one row. Then the elements of matrix looks a graph of 1-dimentional function of 2-dimentional region. When function f(red curve)
is applied to f itself, one can plot \(f(x)\cdot f(x)\) following f(x) value for each x coordinate. one peak function is converted to 2 peak function, 4 peak function ,8 peak \ldots{} and so on.

\begin{figure}
\centering
\includegraphics[width=.9\columnwidth]{img/folding_simple.pdf} 
\label{fig:folding}
\caption{folding mechanism}
\end{figure}

In case \(W\)s are diagonal matrices, consider product of matrix \(Q^T\) ,K and attention map A, if
the shape of elements of Q is same as the ones of K, one can fold graph
of $Q^T$ or K by matrix product. This can be thought the attention matrix A, which is the product of Q and K.

When calculating the row of A where blue point exist, multiply product
of blue row of Q and each column of K. In this case the only elements that are nonzero are those resulting from the dot product with the
column of K (blue) that has nonzero elements in the same number of rows as that number of columns; consequently, only the elements of A that are
in the same row as the blue elements of Q and the same column as the
blue elements of K are nonzero. The self reference structure can be achieved by this operation.

Since \(W_q\) is an identity matrix, it might seem sufficient to simply use K to expand it to an nxn shape, but this approach appears somewhat
passive when compared to the original meanings of ``Key'' and ``Query''.
In other words, one could say that attention performs more self-references in a single iteration. It might also be said that the
adjustment of ε in the function map is performed internally by the model.

\begin{figure}
\centering
\includegraphics[width=.8\columnwidth]{img/attention.pdf}
\label{fig:attention}
\caption{attention matrix calculation}
\end{figure}

The above example is only one element of each rows is non-zero, in case
multiple elements are non-zero FD can be thought as converter of multi
valued function of probablistic distribution function (probablistic
process).In transformers renormalization for probablistic distribution
function is calulated by softmax and division by \(\sqrt{d}\). In this
paper we only treat same weight parameter for each layers compare to FV
and attentions, recently this is called looped transformer and paid
attention with researchers\cite{loopedTr}.

One of the interesting property of FD is hierachical structure of points. Fixed points are on diagonal line called type I, type II fixed
points is depends on type III fixed points refer to \ldots and so on\cite{FD2}. 

\begin{figure}
\centering
\includegraphics[width=.9\columnwidth]{img/hierchical.pdf}
\label{fig:hierachical}
\caption{hierachical structure of points of functional dynamics}
\end{figure}

This hierrachical structure is not merely analogy of the one of
natural/programming languages but coreresponds to the deduction or
in-context learning process of transformers. As in following figure,
functional dynamics can generate self similar fractal shaped function by
adding matrix operation as in attention mechanism. The compsition of
attention (\(f\cdot f\)) and MLP as operators makes self recuesivee
fractal shaped function easily. Fig . shows ssteps to make make two
identical map inside the region of s map. This fact also implies self
similar structure of language related to folding mechanism of FD.

\begin{figure}
\centering
\includegraphics[width=.9\columnwidth]{img/FMAP_INCURSIVE.pdf}
\label{fig:incursive}
\caption{incursive function dynamics using (\(f\cdot f\)) and matrix transpose and multipricalition}
\end{figure}

The original form of FD only consists of function apply(\(\cdot\)),addition (+) and multiplication of constant value \(\epsilon\) this
restrict related to logic structure which transformers can calculete as
following chapter.

\hypertarget{category-theory}{%
\subsection{Category Theory}\label{category-theory}}

As described above,section FD and attention can be treated as some kind
of Category and it should have ability to explain and evaluate
functions. Lambda calculus treats functions as same as variables. All
calculation in is multiple steps of evaluations(eval) and
applications(apply) of formulars. Eval is so called charactor string as
a formular and calucule this, apply is the process that substitutiig
eval's result to other formular. This eval-apply loop is common at the
various field of computer programming. Lambda calcuals has three rules,
alpha conversion beta reduction and eta conversion. Alpha conversion is
just replacement of bound variable names. Beta reduction is application
of a function described by \((\lambda x. f) b=f(b)\) in usual notation.
Eta conversion is desciribed as \(\lambda x. f\) x=f , here rhs and
lhs are same function (constant). This is corresponds to extentionality
definition of functions sets theorem.

Lambda calculus is formulate by using Cartesian closed category (CCC) which have product $X \times Y $ of tow objects X,Y and exponential
object \(X^Y\). There is natual bjiction \(Hom(X \times Y,Z) \simeq Hom(X,Z^Y)\). There is one morphism called
carring \(\lambda g\) for all g and morphism \(Z^Y \rightarrow Z\) is evaluation of
program(S-expression in LISP),this is coresspons to calculation of Attention matrix from Q and K, \(f\cdot f\) of functional dynamics.

\begin{figure}
\centering
\includegraphics[width=.3\columnwidth]{img/eval.pdf}
\label{fig:eval}
\caption{$\lambda$ and eval morphism in CCC}
\end{figure}

There is another least restricted category called Symmetric Monoid Closed Category(SMCC). SMCC do not have diagonal morphism. Intuidively
diagonal morphism and its dual is copy and delete operation. When logic and proof process changes called linear logic. This condition is common
when the objects are vector space and morphisms are linear transformation because \(X\times X=X^2\) is nonlinear. The category called \(\bf{Vect}\).

Be aware with cardinality of exponential object is larger than the cardinality of objects as Lawvere's fixed-point theorem\cite{lawveresfixedpoint} says.

Markov category(MC) is a modeling of probablistic calculation and statistical inference and induction. The object are probablistic
distributions, the morphism are transition kernels between distributions. Generally MC is not CCC,

Topos is defined CCC which has subobject classifiers \(Sub: C^{op}\rightarrow Set\).

Intuitively function f and g of FD corresponds to morphisms, the functor is functional dynamics. In other formulation f,g are objects, functional dynamics itself is morphism and the functor is parametrize by \(\epsilon\). By restricting the formular of FD linear interpolations as in the original paper\cite{FD1}, category theory can explation its parameters \(\epsilon\). Functor between FD and parameter \(\epsilon\)
can be thought natual transformation.

Yoneda's lemma explains the relation between the behavior and parameters FD. According to Yoneda's lemma, for a category C, and a functor from C to othre category D. \(F(x) \leftrightarrow , x \in C, F: C^{op}\)
C\^{}\{op\}\$ is the opposite category of C which morphics is reverse directino to C.

Hom functor \(Hom_C(-,X)\) maps object A to a set of morphisms \(Hom\_C(A,X)\), morphism f: to set of morphism of morphism. There are
natural transformations \(Nat(h_X,h_Y) (h_X \rightarrow h_Y)\). Yoneda's lemma states for any functor F \(Nat(h\_X,F) \simeq F(X) \). This means
a natrual transformation between FD F and G corresponds to a specific
parameter. In case of linear interpolation FD, this is parameter \(epsilon\). Here we identify attention matrix dynamics along layers is
functional dynamics. The change of attention matrix dynamics along to layers can be said natural transformation and there is a specific object(values) maped by a functor F.

\hypertarget{existing-categorical-foundations-of-machine-learning}{%
\subsubsection{Existing categorical foundations of machine learning}\label{existing-categorical-foundations-of-machine-learning}}

Our decomposition builds on an established body of categorical machine-learning theory, which we use rather than re-derive.

\textbf{Gradient-based learning as parametric lenses.} Fong, Spivak and Tuyéras \cite{fong2019backprop} framed supervised learning functorially (``backprop as functor''). Cruttwell, Gavranović, Ghani, Wilson and Zanasi \cite{cruttwell2022categorical} gave a categorical semantics of gradient-based learning in terms of \emph{parametrised maps} (\(\mathrm{Para}\)), \emph{lenses},
and \emph{reverse derivative categories} (RDC), unifying optimisers (SGD, Adam, Nesterov) and losses (MSE, softmax cross-entropy) as
instances of a single parametric-lens structure. Elliott \cite{elliott2018simple} gave
the ``simple essence'' of automatic differentiation as compilation to categories. These works supply the \emph{learning} side of the picture;
they are agnostic to the forward architecture.

\textbf{Reverse derivative categories.} Cockett et al.\cite{CRDC}
axiomatize the reverse derivative \(R[f]:A\times B\to A\) (backpropagation, \(J^\top y\) in \(\mathbf{SMOOTH}\)) and show it
equals a forward derivative plus a dagger structure on the subcategory
of linear maps, which form an additively enriched category with dagger biproducts. This is the setting in which our two Jacobians (state vs parameter) and the additive biproduct structure of.

\textbf{Architectures as (co)algebras.} Categorical Deep Learning\cite{gavranovic2024categorical} present architectures as (co)algebras over (co)monads via polynomial functors. O'Neill et al.\cite{oneill2025selfattentionparametricendofunctorcategorical} specifically formalize self-attention as a parametric endofunctor in \(\mathrm{Para}(\mathbf{Vect})\), restricted to the linear components --- directly the SMCC/value-path fragment we use in \ref{attention-mechanism-transformers-and-their-components}.

\textbf{Probabilistic structure.} Fritz \cite{Markov} axiomatizes Markov categories; the Kleisli category of the distribution monad is the
canonical example. Shiebler, Gavranović and Wilson \cite{shiebler2021category}, surveying category theory in machine learning, already treat learners as
parametrised stochastic processes and distinguish a \emph{co-Kleisli} composition (in which randomness is shared) from a
\emph{\(\mathrm{Para}\)} composition --- a precedent for combining probabilistic (Markov) with parametric structure, which is exactly the
join we make explicit for attention in §6--7.

\textbf{Expressivity via topos theory.} Villani and McBurney \cite{Topos} analyze the \emph{expressivity} of the architecture through topos
theory. They decompose self-attention into a \textbf{choose} morphism --- an input-dependent selection of parameters, which parameterises a
feed-forward architecture --- followed by an \textbf{eval} morphism that computes the output using the chosen architecture; the layer is thus a
linear function whose coefficients are computed by a function that
outputs a function, to which a canonical evaluation is then applied. On
this basis, convolutional, recurrent and graph-convolutional networks
embed in a pretopos of piecewise-linear (PL) functions, whereas the
transformer necessarily lives in its topos completion, so that the two
families instantiate different fragments of logic: the former first-order, the transformer a higher-order reasoner.

This is an important and, we believe, correct result \emph{about expressivity}, and their choice of a PL base is well-motivated: ReLU
networks are exactly piecewise-linear. Their choose-eval decomposition is an independent, architecture-level arrival at the same eval--apply reading we develop from Function Vectors\cite{FV} \ref{the-first-experiment-controlled-probes-on-a-synthetic-in-context-model}: their input-dependent ``choose'' is our extraction \(E:\mathrm{Ctx}\to P_{FV}\), and their canonical evaluation is our counit \(\mathrm{ev}\). We do not contest their expressivity claim; \ref{two-different-quantities-are-being-measured} §9
explains why the present paper nonetheless measures a different quantity. (See also Belfiore and Bennequin \cite{Belfiore} for topos-and-stack models of eep networks.)

\textbf{The linear-logic side.} SMCCs are the standard categorical models of multiplicative intuitionistic linear logic; the SMCC $\leftrightarrow$ (linear) λ-calculus correspondence is the linear analogue of the Curry--Howard--Lambek correspondence between CCCs and the simply-typed λ-calculus, with roots in the differential/linear-logic tradition (Ehrhard--Regnier). This is what licenses the linear-λ reading in \ref{linear-logic-linear-lambda-calculus}.

\textbf{Mechanistic anchors.} Geva et al. \cite{Geva} show feed-forward layers act as key-value memories writing into the residual
stream (§4). Todd et al.\cite{Todd} introduce Function Vectors, compact task representations that \emph{trigger} rather than perform a task, invoking Church's $\lambda$-calculus (§3.3, §8.2).

\hypertarget{linear-logic-linear-lambda-calculus}{%
\subsection{Linear Logic, Linear lambda calculus}\label{linear-logic-linear-lambda-calculus}}

Linear logic is restriction of usual mathematical logic which only allows finite use of propositions during a deduction.

operator \(A \multimap B\) means is linear implication, which signifies ``deriving a conclusion by consuming a premise exactly once''.

\hypertarget{formulation}{%
\section{Formulation}\label{formulation}}

\hypertarget{summary-related-categories}{%
\subsection{Summary related categories}\label{summary-related-categories}}

The claim examined here is narrow and structural: a Transformer layer is not modelled by one symmetric monoidal closed category, but by a \textbf{composite}

\[
\underbrace{\mathrm{Kl}(D)}_{\text{softmax: Markov category}}
\;\xrightarrow{\;D\text{-algebra (expectation)}\;}\;
\underbrace{(\mathbf{Vect},\otimes,\multimap)}_{\text{value path: SMCC} \;\leftrightarrow\; \text{linear }\lambda\text{-calculus}}
\;\xrightarrow{\;\Phi=\text{MLP}\;}\;
\underbrace{(X\multimap Y)}_{\text{realization / apply}}
\;+\;
\underbrace{(\mathbf{Vect},\oplus)}_{\text{residual: additive copy}} .
\]

Four sub-claims make this up: (i) the attention value path is an SMCC
morphism and therefore admits a linear-λ reading (§3); (ii) the MLP is a
realization map / key--value memory (§4); (iii) the residual connection
realizes a copy, but only in the additive sense (§5); (iv) softmax is a
Markov-category morphism (§6). Section 7 assembles the composite, §8
reports experiments, and §9 positions the account against prior work.

\hypertarget{self-attentions-value-path-as-an-smcc-and-the-linear-ux3bb-reading}{%
\subsection{Self-attention's value path as an SMCC, and the linear-λ
reading}\label{self-attentions-value-path-as-an-smcc-and-the-linear-ux3bb-reading}}

\hypertarget{the-smcc-and-its-internal-language}{%
\subsubsection{The SMCC and its internal
language}\label{the-smcc-and-its-internal-language}}

Work in \((\mathbf{Vect}_k,\otimes,I=k)\), a symmetric monoidal closed
category with internal hom \(A\multimap B=\mathrm{Hom}_k(A,B)\),
tensor--hom adjunction
\(\mathrm{Hom}(C\otimes A,B)\cong\mathrm{Hom}(C,A\multimap B)\)
(currying), and evaluation counit

\[\mathrm{ev}_{A,B}:(A\multimap B)\otimes A\longrightarrow B,\]

which is bilinear, hence a genuine linear morphism from the tensor. This
is the categorical content of ``evaluation is available for linear
maps''. Crucially \(\mathbf{Vect}\) has \textbf{no} natural diagonal
\(\Delta:A\to A\otimes A\) (the map \(a\mapsto a\otimes a\) is
non-linear), so copy/delete are absent from the multiplicative fragment.
The internal language of this SMCC is MILL \((\otimes,\multimap,I)\):
types are objects, terms are morphisms, \(\multimap\)-introduction is
currying, \(\multimap\)-elimination is \(\mathrm{ev}\), and the 
linearity discipline (``each variable used exactly once'') is the categorical absence of the diagonal.

\hypertarget{the-linear-approximation}{%
\subsubsection{The linear approximation}\label{the-linear-approximation}}

To keep evaluation non-trivial we freeze the softmax pattern (treat the
attention weights as a query-independent constant) while retaining the
\emph{bilinear} value path \(V=W_v h\) and the query--value contraction.
A fully linearized forward pass would collapse an additive intervention
to a translation and trivialize \(\multimap\); the frozen-pattern
approximation instead preserves the counit \(\mathrm{ev}\). All λ-typing
below is under this approximation; §6 restores the softmax as a Markov
morphism.

\hypertarget{function-vectors-as-points-realization-and-application}{%
\subsubsection{Function Vectors as points; realization and
application}\label{function-vectors-as-points-realization-and-application}}

The reference Function-Vector construction sums, over the top heads, the
out-projection of head activations at the last token, yielding a vector
in \(\mathbb{R}^d=\mathbb{R}^{\mathrm{resid\_dim}}\). An FV is therefore
a \textbf{point} of a parameter object \(P_{FV}\subset\mathbb{R}^d\),
not an element of the internal hom
\(X\multimap Y\cong\mathbb{R}^{d\times d}\). Bridging this dimension gap
requires a \textbf{realization map} \(\Phi:P_{FV}\to(X\multimap Y)\).
The task index \(t\) is data (inferred from context), not a free
parameter; we separate extraction and application:

\[E:\mathrm{Ctx}\multimap P_{FV},\qquad A:=\mathrm{ev}\circ(\Phi\otimes\mathrm{id}_X):P_{FV}\otimes X\to Y,\qquad v_t=E(\mathrm{ctx}_t).\]

\hypertarget{in-context-computation-as-a-linear-ux3bb-term}{%
\subsubsection{In-context computation as a linear
λ-term}\label{in-context-computation-as-a-linear-ux3bb-term}}

With \(E,\Phi\) as constants and using currying, one forward pass of ICL
is well-typed in linear λ-calculus:

\[
\dfrac{c:\mathrm{Ctx}\vdash \Phi(E\,c):X\multimap Y \qquad x:X\vdash x:X}
{c:\mathrm{Ctx},\,x:X\vdash (\Phi(E\,c))\,x:Y}\;(\multimap\text{-elim}=\mathrm{ev})
\qquad\Longrightarrow\qquad
\boxed{\;\lambda c.\,\lambda x.\,(\Phi(E\,c))\,x\;:\;\mathrm{Ctx}\multimap(X\multimap Y)\;}
\]

Both \(c\) and \(x\) are used exactly once, so the term is linear. The
type \(\mathrm{Ctx}\multimap(X\multimap Y)\) says the context compiles
to a \emph{function-typed value}: the FV \(E(c)\) is a reified function
reference (a first-class datum), \(\Phi\) resolves it to a procedure,
and \(\mathrm{ev}\) applies it --- a precise reading of ``the vector
triggers, but does not perform, the task'' {[}Todd et al.~2024{]}.

\hypertarget{the-residual-stream-as-a-polynomial-extension-the-abstractionapplication-adjunction}{%
\subsubsection{The residual stream as a polynomial extension; the
abstraction--application
adjunction}\label{the-residual-stream-as-a-polynomial-extension-the-abstractionapplication-adjunction}}

The categorical basis of λ-abstraction is Lambek's \textbf{functional
completeness} {[}Lambek 1974{]}, realized through the \textbf{polynomial
extension} \(\mathcal{C}[x]\): adjoining an indeterminate arrow
\(x:1\to A\) yields a category \(\mathcal{C}[x]\) in which every
polynomial morphism corresponds uniquely to a morphism of
\(\mathcal{C}\), so \(\mathcal{C}[x](B,C)\cong\mathcal{C}(A\times B,C)\)
and, by closedness, \(\cong\mathcal{C}(B,A\Rightarrow C)\) ---
categorical λ-abstraction.

Došen {[}2001{]} refines this into two adjunctions. The inclusion
\(I:\mathcal{C}\to\mathcal{C}[x]\) of a cartesian closed category into
its polynomial extension has a \textbf{left adjoint based on product}
(\(A\times-\)) and a \textbf{right adjoint based on exponentiation}
(\(A\Rightarrow-\)); composing the two --- which, as Došen notes, stem
from the deduction theorem --- produces the product--exponentiation
adjunction \(A\times-\dashv A\Rightarrow-\). In it, \textbf{application
is the counit} \(\mathrm{ev}:(A\Rightarrow C)\times A\to C\) and
\textbf{abstraction is the unit}, with β- and η-conversion the two
triangle identities (the inversion principle Došen draws in parallel
with comprehension/extensionality). This \emph{is} the eval--apply
structure: apply \(=\) counit (evaluation), abstract \(=\) unit.

Two adaptations are forced here.

\textbf{(i) Linear polynomial extension.} In the SMCC
\((\mathbf{Vect},\otimes,\multimap)\) the indeterminate must be
\textbf{linear} (used exactly once), so the product-based left adjoint
becomes the tensor \(A\otimes-\) and the exponentiation-based right
adjoint becomes \(A\multimap-\). Došen's composed adjunction becomes the
\textbf{tensor--hom adjunction \(A\otimes-\dashv A\multimap-\)} of §3.1,
whose counit is exactly the evaluation
\(\mathrm{ev}:(A\multimap B)\otimes A\to B\) used to type application in
§3.3--3.4. The ``apply'' of the eval--apply loop is thus not posited but
\textbf{derived}: it is the counit of the adjunction that the (linear)
deduction theorem produces from the polynomial extension.

\textbf{(ii) The residual stream as \(\mathcal{K}[x]\).} Identify the
running residual value with the indeterminate: the residual stream
\emph{is} the polynomial extension \(\mathcal{K}[x]\), threading \(x\)
across depth. By §5 the residual copies \(x\) only \textbf{additively}
(the \(\oplus\)-diagonal), never multiplicatively, so the indeterminate
is genuinely linear --- carried across layers but never squared in a
\(\otimes\)-context. The residual therefore instantiates the linear
polynomial extension whose adjunction yields apply, and the
FV-extraction map \(E\) of §3.3 is the (linear) \textbf{abstraction}
direction --- the unit side --- reifying the context as a function-typed
value.

\textbf{Precision: value level vs type level.} The adjunction furnishes
\emph{linear} abstraction/application: substitution of a once-used
indeterminate, the analogue of a single β-step. Full λ-abstraction with
unrestricted binding and reuse of the bound variable would require the
exponential \(!\), which §5 shows is absent. So \(\mathcal{K}[x]\) and
the Došen adjunction supply application/abstraction at the \textbf{value
level}, not unrestricted type-level λ-abstraction with free duplication
--- consistent with the once-only discipline of the whole account: the
eval--apply loop is linear, and the polynomial extension generating its
apply-counit is a linear one.

\hypertarget{the-mlp-as-realization-and-keyvalue-memory}{%
\subsection{The MLP as realization and key--value
memory}\label{the-mlp-as-realization-and-keyvalue-memory}}

The MLP is best read not as evaluation but as the (non-linear)
realization map \(\Phi:P_{FV}\to(X\multimap Y)\) that converts an FV
point into an applied function. This is consistent with the mechanistic
result that feed-forward layers act as \textbf{key--value memories}:
each key correlates with input patterns and each value induces an output
distribution, with the layer output a composition of memories refined
through the residual stream {[}Geva et al.~2021, 2022{]}. In our
composite the MLP writes additively into the residual (the \(\oplus\)
side of §5) while realizing the currying step of §3. It is therefore the
natural carrier of \(\Phi\), and we test this localization in §8 (Probe
B, Probe C).

\hypertarget{the-residual-connection-as-additive-copy}{%
\subsection{The residual connection as additive
copy}\label{the-residual-connection-as-additive-copy}}

A residual connection is \(x\mapsto x+f(x)=(\mathrm{id}_A+f)(x)\), a sum
of parallel morphisms, legitimate because \(\mathbf{Vect}\) is
additively enriched with biproducts. It uses the
\textbf{additive/biproduct diagonal} \(\Delta^{\oplus}:A\to A\oplus A\),
\(a\mapsto(a,a)\) (fan-out to sublayers) and codiagonal
\(\nabla^{\oplus}:A\oplus A\to A\) (write-back), both linear. Three
notions of copy must be separated:

\begin{longtable}[]{@{}lllll@{}}
\toprule
\begin{minipage}[b]{0.17\columnwidth}\raggedright
Notion\strut
\end{minipage} & \begin{minipage}[b]{0.17\columnwidth}\raggedright
Structure\strut
\end{minipage} & \begin{minipage}[b]{0.17\columnwidth}\raggedright
Axis\strut
\end{minipage} & \begin{minipage}[b]{0.17\columnwidth}\raggedright
Carrier\strut
\end{minipage} & \begin{minipage}[b]{0.17\columnwidth}\raggedright
Status\strut
\end{minipage}\tabularnewline
\midrule
\endhead
\begin{minipage}[t]{0.17\columnwidth}\raggedright
Multiplicative diagonal\strut
\end{minipage} & \begin{minipage}[t]{0.17\columnwidth}\raggedright
\(\Delta:A\to A\otimes A\)\strut
\end{minipage} & \begin{minipage}[t]{0.17\columnwidth}\raggedright
---\strut
\end{minipage} & \begin{minipage}[t]{0.17\columnwidth}\raggedright
none\strut
\end{minipage} & \begin{minipage}[t]{0.17\columnwidth}\raggedright
Forbidden (linearity, ``use once'')\strut
\end{minipage}\tabularnewline
\begin{minipage}[t]{0.17\columnwidth}\raggedright
Additive diagonal\strut
\end{minipage} & \begin{minipage}[t]{0.17\columnwidth}\raggedright
\(\Delta^{\oplus}:A\to A\oplus A\)\strut
\end{minipage} & \begin{minipage}[t]{0.17\columnwidth}\raggedright
depth (across layers)\strut
\end{minipage} & \begin{minipage}[t]{0.17\columnwidth}\raggedright
residual stream\strut
\end{minipage} & \begin{minipage}[t]{0.17\columnwidth}\raggedright
Legal, linear\strut
\end{minipage}\tabularnewline
\begin{minipage}[t]{0.17\columnwidth}\raggedright
Markov copy\strut
\end{minipage} & \begin{minipage}[t]{0.17\columnwidth}\raggedright
non-natural comonoid in \(\mathrm{Kl}(D)\)\strut
\end{minipage} & \begin{minipage}[t]{0.17\columnwidth}\raggedright
position (across tokens)\strut
\end{minipage} & \begin{minipage}[t]{0.17\columnwidth}\raggedright
softmax\strut
\end{minipage} & \begin{minipage}[t]{0.17\columnwidth}\raggedright
Legal, generates correlation\strut
\end{minipage}\tabularnewline
\bottomrule
\end{longtable}

So the residual connection \emph{does} recover a copy --- the
\textbf{additive} one, along the depth axis --- which is why ``copy
exists in some sense'' is correct. But it is not the multiplicative
diagonal, and it is not the exponential modality \(!\): the two branches
are summed back by \(\nabla^{\oplus}\) into a single
\(\otimes\)-resource rather than being consumed independently, so the
once-only discipline of \(\mathrm{ev}\) is untouched. The residual
therefore adds the \emph{additive} fragment \((\oplus,\&)\) to the
multiplicative \((\otimes,\multimap)\) without breaking linearity, and
does not supply unbounded reuse. In the polynomial-extension reading of
§3.5, this is exactly why the residual realizes a \emph{linear}
\(\mathcal{K}[x]\): the additive diagonal carries the indeterminate
\(x\) across depth, but the absence of the multiplicative diagonal keeps
each occurrence linear, so the abstraction--application adjunction it
supports is the linear (value-level) one, not the \(!\)-enabled type-level one.

\hypertarget{softmax-as-a-markov-category-morphism}{%
\subsection{Softmax as a Markov-category morphism}\label{softmax-as-a-markov-category-morphism}}

Each row of \(A=\mathrm{softmax}(QK^\top)\) is a probability
distribution over key positions. Hence \(A\) is a \textbf{Markov
kernel}: a morphism of the Kleisli category \(\mathrm{Kl}(D)\) of the
distribution monad \(D\) (equivalently, \(\mathrm{FinStoch}\)), the
canonical Markov category {[}Fritz 2020{]}. The monad \(D\) is the
\emph{generator} of this Markov category --- \(\mathrm{Kl}(D)\) is a
Markov category because \(D\) is a commutative affine monad
(\(D(1)\cong 1\) yields the unique discarding map; commutativity yields
the symmetric monoidal structure). Value mixing \(A\,V\) is the action
of the associated \textbf{\(D\)-algebra} \(\mathbb{E}:D(V)\to V\)
(convex combination / expectation), which aligns with the Kolmogorov
picture in which random variables are functions on a sample space and
expectation is a linear functional. A head thus decomposes as

\[\text{head}=\Big[\;\mathrm{Ctx}\xrightarrow{\ \text{softmax}\ } D(X)\;\Big]\ \text{followed by}\ \Big[\;D(X)\xrightarrow{\ \mathbb{E}\ } X\;\Big],\]

a Markov kernel composed with the \(D\)-algebra structure map --- the
probabilistic generalization of the SMCC evaluation of §3. The Markov
copy is \emph{non-natural}: copy-then-randomize differs from randomize-then-copy, so it generates correlation rather than duplicating
a value freely; this is the position-axis copy of §5 and, as we note in §7, it is not the exponential \(!\).

\hypertarget{the-composite-markov-category-smcc-additivemlp}{%
\subsection{The composite: Markov category $\cdot$ SMCC (+ additive/MLP)}\label{the-composite-markov-category-smcc-additivemlp}}

Assembling §§3--6, a Transformer layer is the composite displayed in §1:
a Markov morphism (softmax) feeding, through the \(D\)-algebra expectation, the SMCC value path (where evaluation and the linear-$lambda$ reading live), then the non-linear realization \(\Phi\) (MLP), with the
additive residual \((\mathbf{Vect},\oplus)\) threading depth. Three
points make the composite non-collapsible into a single closed category.

\begin{enumerate}
\def\labelenumi{\arabic{enumi}.}
%%\tightlist
\item
  \textbf{The closed structure comes from the SMCC side, not the Markov
  side.} Markov categories are generally not closed; the internal hom
  that supports \(\mathrm{ev}\)/apply is \(\mathbf{Vect}\)'s
  \(\multimap\). A single ``closed Markov category'' cannot host both
  the non-natural copy and the closed structure; the correct object is
  the \emph{composite} of two categories bridged by the \(D\)-algebra.
\item
  \textbf{Copy differs on each side.} The value/eval SMCC has \emph{no}
  copy (once-only); the residual has \emph{additive} copy; softmax has
  \emph{non-natural Markov} copy. None is the exponential \(!\), so the
  whole layer has no unbounded-reuse capacity.
\item
  \textbf{The learning side is orthogonal and already categorified.}
  Parameter updates are handled by the \(\mathrm{Para}\)/lens/RDC
  machinery of the existing foundations {[}Cruttwell et al.~2022;
  Cockett et al.~2020{]}; our contribution is the \emph{forward-architecture} decomposition, extending the linear-endofunctor picture of O'Neill et al.~{[}2025{]} with the Markov (softmax) and additive (residual) structure, in the spirit of
  the Kleisli+\(\mathrm{Para}\) combination already anticipated by   Shiebler et al.~{[}2021{]}.
\end{enumerate}

\hypertarget{experiments}{%
\section{Experiments}\label{experiments}}

\hypertarget{the-first-experiment-controlled-probes-on-a-synthetic-in-context-model}{%
\subsection{The first experiment, Controlled probes on a synthetic in-context model}\label{the-first-experiment-controlled-probes-on-a-synthetic-in-context-model}}

We train a small decoder Transformer (\(V=10\) symbols, \(k=5\) in-context examples, \(d=64\), \(L=3\) layers, \(4\) heads) on an
in-context permutation-application task (infer a random bijection \(\pi\) from examples, apply it to queries). This is a methodology
demonstrator, not a language model; the same probes attach to real models via forward hooks. The model reaches \textbf{single-query accuracy \(1.000\)} (chance \(0.100\)).

\textbf{Probe A --- copy localization (sublayer-local cross-position
mixed second difference).} Perturbing a sublayer's own input at two
distinct source positions \(a\neq b\) and reading
\(f(+a,+b)-f(+a)-f(+b)+f()\) at the query position isolates each
sublayer's own cross-position (\(\otimes\)/copy) interaction; a strictly
position-wise MLP generates none of its own.

\begin{longtable}[]{@{}llll@{}}
\toprule
Layer & attention & MLP & ratio attn/mlp\tabularnewline
\midrule
\endhead
0 & \(5.47\times10^{-3}\) & \(0.0\) &
\(\sim 5\times10^{6}\)\tabularnewline
1 & \(4.11\times10^{-4}\) & \(0.0\) &
\(\sim 4\times10^{5}\)\tabularnewline
2 & \(5.07\times10^{-4}\) & \(0.0\) &
\(\sim 5\times10^{5}\)\tabularnewline
\bottomrule
\end{longtable}

The MLP cross-position term is \textbf{machine zero} at every layer;
only attention carries cross-position interaction. This localizes the
copy/diagonal structure in softmax attention, not in the pointwise MLP
nonlinearity --- the empirical counterpart of §5--6.

\textbf{Probe B --- apply localization (activation patching,
flip-to-clean rate).} Patching the clean run's activation at the query
position into a corrupted run (different bijection, same query) shows
which component transfers the applied answer (356 usable trials).

\begin{longtable}[]{@{}llll@{}}
\toprule
Component & Layer 0 & Layer 1 & Layer 2\tabularnewline
\midrule
\endhead
attention-out & \(0.00\) & \(\mathbf{1.00}\) & \(0.00\)\tabularnewline
MLP-out & \(0.00\) & \(0.32\) & \(\mathbf{1.00}\)\tabularnewline
\bottomrule
\end{longtable}

Mid-layer attention (L1) fully transfers the answer, and the final-layer
MLP (L2) fully writes it; last-layer attention does not. This supports
\textbf{evaluation/routing in attention} and \textbf{realization
\(\Phi\) in the final MLP}, matching §3--4.

\textbf{Probe C --- reuse ablation (per-query accuracy on multi-query
prompts).}

\begin{longtable}[]{@{}llll@{}}
\toprule
Condition & q0 & q1 & q2\tabularnewline
\midrule
\endhead
intact & \(1.000\) & \(1.000\) & \(1.000\)\tabularnewline
freeze attention (uniform) & \(0.047\) & \(0.027\) &
\(0.090\)\tabularnewline
linearize MLP (remove nonlinearity) & \(0.148\) & \(0.166\) &
\(0.180\)\tabularnewline
\bottomrule
\end{longtable}

Freezing attention collapses multi-query reuse to chance (the Markov
copy of §6 is what routes one function to many arguments), while
linearizing the MLP degrades the applied value. Both components are
necessary, in the two distinct roles predicted: attention carries
routing/copy, MLP carries the applied value.

Together the probes support \textbf{eval = attention (with the copy), apply/realization = MLP}, and reject either monolithic reading.


\hypertarget{the-second-experiment-architecture-comparison}{%
\subsection{The second experiment, architecture comparison}\label{the-second-experiment-architecture-comparison}}

Another question is that wehre outstanding performance of transformers
comes from and is it related to category theoretical sturucture shonw in
this paper or not. There is another possibility that architecture of
network is not main cause of performance, specific learned weight values
are essential for the performance. To solve this problem, we prepare an
experiment comparing accuracy of retrival task with recurrent neural
network (RNN) and state space model (SSM), which transformers good at
and tatget tasks of function vector study{[}FV{]}.

\hypertarget{three-architecture-comparisonm4chance0.167}{%
\subsubsection{three architecture
comparison(m=4、chance=0.167)}\label{three-architecture-comparisonm4chance0.167}}

\begin{longtable}[]{@{}lll@{}}
\toprule
Transformer & RNN(GRU) & diagonal SSM\tabularnewline
\midrule
\endhead
\textbf{1.000} & 0.405 & 0.322\tabularnewline
\bottomrule
\end{longtable}

In the Accurary retrieval task, the Transformer handily defeated
state-based models. The difference was evident between the
Transformer---which can address arbitrary past constraints---and
RNNs/SSMs, which compress context into a fixed-size state. Since we
evaluated retrieval (which the Transformer excels at) rather than state
tracking (which the Transformer struggles with).

\hypertarget{components-ablationtransformerm4}{%
\subsubsection{components ablation(Transformer、m=4)}\label{components-ablationtransformerm4}}

\begin{longtable}[]{@{}lll@{}}
\toprule
intact & freeze\_attn(freeze dynamic routing) & linear\_mlp(freex eval)\tabularnewline
\midrule
\endhead
1.000 & \textbf{0.284} & 0.855\tabularnewline
\bottomrule
\end{longtable}

In ablation testing to isolate which components are responsible for this
performance, a clear distinction in accuracy emerged. For m=4,
Transformer=1.000, RNN=0.405, and SSM=0.322 (chance 0.167)---in the
retrieval experiment, Transformer decisively outperformed the static
model.

The theory presented in this paper predicts that ``routing =
softmax(Markov kernel), apply = eval,'' and retrieval is precisely the
task that performs half of that routing. The fact that this broke down
with frozen attention serves as evidence that the architecture predicted
by the theory is indeed being utilized in retrieval; it is not a failure
of the categorical framework but rather a success of the architecture.
This confirms the assertion that ``Transformer's success in retrieval
stems from the success of the Markov (softmax = dynamic routing) + SMCC
(eval) architecture, not a failure of that framework.''

\emph{Expected reading:} attention-dominant cross-position term corroborates the softmax (Markov) copy of §6, with the toy model providing the exact position-wise isolation.

\emph{Epistemic status.} these experients are constructions under an explicit linear approximation (frozen softmax for the λ-typing; the Markov reading of §6 restores it),and  reports a small synthetic model as a methodology demonstrator, not evidence at language-model scale.

\hypertarget{related-work-and-positioning-expressivity-versus-resource-structure}{%
\section{Related work and positioning: expressivity versus resource
structure}\label{related-work-and-positioning-expressivity-versus-resource-structure}}

The \emph{learning} side of Transformers is already on firm categorical footing: parametric lenses and reverse derivative categories
\cite{CRDC}\cite{fong2019backprop}\cite{cockett2020reverse,}, surveyed in \cite{shiebler2021category}, and the (co)algebraic architecture theory of \cite{gavranovic2024categorical}. The \emph{forward} structure of attention was categorified for its linear
part by O'Neill et al.\cite{oneill2025selfattentionparametricendofunctorcategorical}. Our account adds the two structures that the linear picture defers: the \textbf{Markov} structure of softmax \cite{Markov} and the \textbf{additive} structure of residuals, and
states how they compose with the SMCC value path and the MLP
realization. The combination of a probabilistic (Kleisli) and a parametric structure was already anticipated in the survey of Shiebler et al.~\cite{shiebler2021category}; we make it concrete for attention and connect it, via the linear-λ reading, to the Function-Vector phenomenology of \cite{FV} and the key--value-memory role of the MLP \cite{geva2021transformerfeedforwardlayerskeyvalue}.

\hypertarget{two-different-quantities-are-being-measured}{%
\subsection{Two different quantities are being measured}\label{two-different-quantities-are-being-measured}}

The relation to the topos analysis of Villani and McBurney \cite{Topos} deserves care, because a superficial reading makes our accounts look contradictory: they place the transformer in a \textbf{topos} (cartesian closed, with free copy and a full higher-order internal logic), while we
place the value path in a \textbf{SMCC} (monoidal closed, no multiplicative diagonal, hence \emph{linear} λ-calculus). We stress that
this is \textbf{not a disagreement about a shared question, and we allege no flaw in their argument}. Given their setting --- a PL base
category, which is the right home for ReLU networks --- their conclusion follows. Cartesian closure yields greater expressive power than linear
logic; linear logic is not weaker in what it can \emph{compute} (it embeds intuitionistic logic via \(A\to B \;=\; {!A}\multimap B\)) but
\emph{finer} in what it can \emph{track}.

The two frameworks measure different quantities:

\begin{longtable}[]{@{}lll@{}}
\toprule
\begin{minipage}[b]{0.30\columnwidth}\raggedright
\strut
\end{minipage} & \begin{minipage}[b]{0.30\columnwidth}\raggedright
Villani--McBurney {[}2024{]}\strut
\end{minipage} & \begin{minipage}[b]{0.30\columnwidth}\raggedright
This paper\strut
\end{minipage}\tabularnewline
\midrule
\endhead
\begin{minipage}[t]{0.30\columnwidth}\raggedright
Question\strut
\end{minipage} & \begin{minipage}[t]{0.30\columnwidth}\raggedright
What can the architecture \textbf{compute}? (expressivity)\strut
\end{minipage} & \begin{minipage}[t]{0.30\columnwidth}\raggedright
How often is each value \textbf{used} in one forward pass? (resource
structure)\strut
\end{minipage}\tabularnewline
\begin{minipage}[t]{0.30\columnwidth}\raggedright
Base category\strut
\end{minipage} & \begin{minipage}[t]{0.30\columnwidth}\raggedright
PL functions → topos completion\strut
\end{minipage} & \begin{minipage}[t]{0.30\columnwidth}\raggedright
\((\mathbf{Vect},\otimes,\multimap)\) + \(\mathrm{Kl}(D)\) +
\((\mathbf{Vect},\oplus)\)\strut
\end{minipage}\tabularnewline
\begin{minipage}[t]{0.30\columnwidth}\raggedright
Logic\strut
\end{minipage} & \begin{minipage}[t]{0.30\columnwidth}\raggedright
Higher-order, copy free\strut
\end{minipage} & \begin{minipage}[t]{0.30\columnwidth}\raggedright
\(!\)-free intuitionistic linear logic\strut
\end{minipage}\tabularnewline
\begin{minipage}[t]{0.30\columnwidth}\raggedright
Verdict\strut
\end{minipage} & \begin{minipage}[t]{0.30\columnwidth}\raggedright
Transformer is a higher-order reasoner\strut
\end{minipage} & \begin{minipage}[t]{0.30\columnwidth}\raggedright
Value application is resource-linear\strut
\end{minipage}\tabularnewline
\bottomrule
\end{longtable}

Expressivity is not our subject, and on that axis we defer to them. A
cartesian model, however, cannot \emph{see} the quantity we are after:
in a category where copy is free, the statement ``this value is consumed
exactly once'' is not expressible. That is precisely the statement we
want to test, because it is what determines whether an in-context
function can be \emph{reused} without being recomputed.

\hypertarget{what-the-pltopos-scope-leaves-out}{%
\subsection{What the PL/topos scope leaves
out}\label{what-the-pltopos-scope-leaves-out}}

Two layers become invisible from a PL base --- by scope, not by error.

First, \textbf{softmax is not piecewise linear.} It contains an
exponential and a normalisation, so it is neither PL nor a composite of
PL maps, and it is the very component that makes ``choose''
input-dependent. A PL account must either freeze it, approximate it, or
abstract it into the choose morphism; in all three cases its
probabilistic structure --- the Markov kernel of §6 --- is not
represented. ReLU MLPs \emph{are} exactly PL, so the PL base is apt for
the MLP; it is the attention routing that escapes it.

Second, \textbf{the resource discipline of application is invisible in a
cartesian setting}, as noted above.

\hypertarget{the-resulting-claim-and-its-limits}{%
\subsection{The resulting claim, and its
limits}\label{the-resulting-claim-and-its-limits}}

Within the linear approximation of §3.2, the value path is a morphism of
a SMCC, whose internal language is the multiplicative fragment
\((\otimes,\multimap)\). The residual connection adds the
\textbf{additive} fragment \((\oplus,\&)\) --- it does \emph{not} relax
the resource discipline, because the two branches of \(\Delta^{\oplus}\)
are summed back by \(\nabla^{\oplus}\) into a single
\(\otimes\)-resource, so no independent second consumption is created
and the exponential \(!\) is not supplied (§5). The correct description
is therefore \emph{addition of the additive fragment}, not
\emph{relaxation of linearity}.

Consequently we can state, for the linear/PL-input regime:

\begin{quote}
\textbf{Claim.} Under the linear approximation, the computation
performed by the value path of a Transformer --- including the copy
afforded by residual connections --- is expressible in
\textbf{\(!\)-free intuitionistic linear λ-calculus}
\((\otimes,\multimap,\oplus,\&)\). Nothing in this regime requires the
unbounded reuse that a full higher-order λ-calculus would license, and
the hypothesis that the model performs only such \(!\)-free linear
computation \textbf{cannot be ruled out}.
\end{quote}

This is not merely a modal claim about what cannot be excluded; §8
supplies positive evidence for it. The multiplicative (cross-position)
interaction that a genuine \(\otimes\)-diagonal would require is
\textbf{absent from the MLP} (Probe A: machine zero at every layer) and
present only in attention, and multi-query reuse \textbf{collapses to
chance when attention is frozen} (Probe C), while surviving MLP
linearization in part. If the model were exploiting a \(!\)-like
duplication implemented in the pointwise nonlinearity, neither result
would hold. The reuse we observe is therefore better attributed to the
\emph{non-natural Markov copy} of softmax (§6) than to an exponential
modality.

\textbf{Limits.} The claim is scoped to the linear approximation with
frozen softmax. Once the softmax is restored, its non-natural Markov
copy is a \emph{candidate} source of duplication, and whether it
effectively supplies something with the strength of \(!\) is open (§10).
Our claim is also about the value path of one layer; it says nothing
about what a deep stack can express in the sense of Villani and
McBurney, where we defer to their result.

\hypertarget{state-tracking-and-the-parallelism-tradeoff-reading-a-three-architecture-experiment-through-complexity-classes}{%
\subsection{State tracking and the parallelism tradeoff: reading a three-architecture experiment through complexity classes}\label{state-tracking-and-the-parallelism-tradeoff-reading-a-three-architecture-experiment-through-complexity-classes}}

The resource-structure axis has an empirical counterpart on the
\emph{hard} side of the architecture --- the tasks a Transformer's
forward category cannot express --- and here our account meets the
circuit-complexity classification of Merrill and collaborators\cite{merrill2023parallelism}. They prove that log-precision Transformers can be
simulated by constant-depth uniform threshold circuits and hence lie in
\(\mathsf{TC}^0\); they frame this as a \emph{parallelism tradeoff} ---
any architecture as parallelizable as the Transformer inherits the same
ceiling. Merrill, Petty and Sabharwal \cite{merrill2024illusion} extend the ceiling to
state-space models: despite their recurrent form, S4- and Mamba-style SSMs are also confined to \(\mathsf{TC}^0\) and, like Transformers,
cannot express permutation composition (\(S_5\)), the canonical inherently-sequential state-tracking problem that even a simple RNN expresses naturally.

We ran a three-architecture capacity sweep on exactly this task ---
explicit composition over \(S_5\), scored on the final state, each
architecture trained under a shared curriculum to remove the
trainability confound --- measuring the largest composition length
\(n^*\) each model masters as a function of its capacity axis
(Transformer heads at fixed head-dimension; RNN and SSM hidden width).
The three architectures separate into three qualitatively distinct, capacity-\emph{independent} plateaus:

\begin{itemize}
%%\tightlist
\item
  \textbf{RNN}: \(n^*\) high and flat (solves the whole ladder at every width) --- the true sequential state update, reaching
  \(\mathsf{NC}^1\);
\item
  \textbf{diagonal SSM}: \(n^*\) at the floor (fails already at short compositions), independent of width --- the empirical face of the
  \emph{Illusion of State} result;
\item
  \textbf{Transformer}: \(n^*\) \emph{intermediate} and flat --- above the SSM, below the RNN, and unmoved by adding heads.
\end{itemize}

Two readings follow, aligned with Merrill et al.'s motivation. First,
the flatness is the parallelism tradeoff seen from the inside: because
the \(\mathsf{TC}^0\) ceiling is a structural consequence of constant
depth and log precision, it is \emph{not} liftable by adding capacity
within the same architecture family. This matches our separate finding
that the Transformer's state-tracking boundary does not scale with depth
either --- the log-graded reading \(n^* \sim c\cdot 2^L\) is not
supported; \(n^*\) is a fixed structural constant. Second, and this is
where our framework adds to theirs, the \emph{heights} of the three
plateaus are explained by the forward category of each architecture.
Merrill et al.~characterize the ceiling from the outside (a circuit
upper bound on what the class can express); we decompose the route to it
from the inside. The ordering RNN \(>\) Transformer \(>\) SSM is the
ordering of copy structure: the RNN's sequential coalgebra supplies a
genuine exponential (unrestricted reuse); the diagonal SSM's
input-independent transition supplies no cross-position coupling at all;
the Transformer's softmax supplies cross-position coupling but only in
the \emph{parallel}, non-natural Markov form of §6, which is bounded
rather than sequential.

Our train-with-ablation experiment localizes the Transformer's
intermediate ability to specific components: removing the softmax (the
Markov kernel) \emph{or} the MLP nonlinearity (the cartesian part) each
drops \(n^*\) to the SSM floor, whereas linearizing the value path
leaves it unchanged. The Transformer's foothold on state tracking ---
its position strictly above the SSM within the shared \(\mathsf{TC}^0\)
ceiling --- thus requires \emph{both} the Markov kernel and the
cartesian nonlinearity together; neither alone suffices. This is the
resource-structure reading of why the Transformer sits where it does in
the complexity classification: it has the cross-position coupling the
SSM lacks, but only the bounded, parallel kind, never the RNN's
sequential exponential.

\textbf{A caveat on registers.} Merrill et al.'s results are asymptotic,
worst-case \emph{expressivity} bounds; our \(n^*\) are finite, learned,
task-specific boundaries. They capture the same phenomenon --- the
\(\mathsf{TC}^0\) ceiling --- from different registers, and should not
be conflated: the Transformer's plateau at a finite \(n^*\) does not
assert that no larger composition is ever solvable, only that within
this scale and training the boundary is capacity-independent. Bridging
the worst-case upper bound and the learned boundary is itself an open direction.

\hypertarget{making-the-composite-cartesian-closed-diagonals-the-lnl-extension-and-external-memory}{%
\section{Making the composite cartesian closed: diagonals, the LNL extension, and external memory}\label{making-the-composite-cartesian-closed-diagonals-the-lnl-extension-and-external-memory}}

A recurring question is whether the composite of §7 --- a Markov
category and an SMCC, plus additive and cartesian fragments --- can be
turned into a \emph{cartesian closed category} (CCC) by supplying the
diagonal \(\Delta_X : X \to X\times X\) that the multiplicative
structure lacks. This section collects the answer, its consequences 
a would-be architectural extension, and the relation to Turing-completeness and chain-of-thought.

\hypertarget{four-ways-to-add-a-diagonal-and-why-each-has-a-cost}{%
\subsection{Four ways to add a diagonal, and why each has a cost}\label{four-ways-to-add-a-diagonal-and-why-each-has-a-cost}}

A CCC is exactly an SMCC together with the exponential modality \(!\) (Seely; Benton's LNL). One cannot make an SMCC cartesian closed without paying somewhere, and the four available routes each surrender a different property.

\begin{enumerate}
\def\labelenumi{\arabic{enumi}.}
%%\tightlist
\item
  \textbf{Add the exponential \(!\) (the linear-logic route).} The
  co-Eilenberg--Moore category of \(!\)-coalgebras is cartesian (Seely),
  so a genuine diagonal lives there. Cost: this is not a description of
  the model but an \emph{extension} --- §5 shows the Transformer has no
  \(!\) (residual copy is additive, softmax copy is non-natural Markov; neither is exponential).
\item
  \textbf{Use the biproduct diagonal \(\Delta_\oplus\) (residual).}
  \((\mathbf{Vect},\oplus)\) is cartesian (§5). Cost: it is not
  \emph{closed} --- \(\oplus\) has no internal hom (§ below), so one
  gets a cartesian but not a cartesian \emph{closed} category; the eval
  lives on \(\otimes\), which has no diagonal.
\item
  \textbf{Restrict to a basis / \(\mathbf{FinSet}\).}
  \(\mathbf{FinSet}\) is a CCC and one-hot vectors admit a diagonal
  \(e_i \mapsto e_i\otimes e_i\). Cost: this destroys linearity (the
  diagonal is not bilinear:
  \(\sum_i c_i e_i \mapsto \sum_i c_i\, e_i\otimes e_i \ne (\sum c_i e_i)\otimes(\sum c_i e_i)\))
  --- i.e.~it collapses the probabilistic superposition that is
  softmax's essence.
\item
  \textbf{Read the Markov copy as a diagonal.} \(\mathrm{Kl}(D)\)
  already has copy (§6, a CD-category). Cost: it is \emph{non-natural}
  --- cartesian diagonals must be natural (commute with every morphism),
  which the Markov copy does only on deterministic (Dirac) morphisms,
  i.e.~after killing the probabilistic content.
\end{enumerate}

All four routes hit the same wall: \textbf{the diagonal (duplication)
and linearity (resource conservation) are incompatible over
probabilistic superposition.} This is not incidental but the structure
theorem CCC \(=\) SMCC \(+\,!\): cartesian closure requires \(!\), and
\(!\) \emph{is} unrestricted duplication, whose introduction abandons
the resource-linearity that is this paper's subject. The three notions
of copy catalogued in §5--6 are, read this way, precisely the three ways
the model \emph{fails to be} cartesian.

\hypertarget{why-mathbfvect-is-not-closed-for-oplus}{%
\subsection{\texorpdfstring{Why \(\mathbf{Vect}\) is not closed for
\(\oplus\)}{Why \textbackslash mathbf\{Vect\} is not closed for \textbackslash oplus}}\label{why-mathbfvect-is-not-closed-for-oplus}}

The obstruction in route 2 deserves a line, since it explains why the
residual's diagonal cannot host an eval. Closure for \(\oplus\) would
require an internal hom \([A,-]_\oplus\) with
\(\mathrm{Hom}(A\oplus B, C) \cong \mathrm{Hom}(B,[A,C]_\oplus)\). But
by the biproduct's universal property
\(\mathrm{Hom}(A\oplus B, C) \cong \mathrm{Hom}(A,C)\times\mathrm{Hom}(B,C)\):
a map out of \(A\oplus B\) is just a \emph{pair} of maps, with nothing
to curry. Dimension-counting makes it sharp:
\(\dim\mathrm{Hom}(A\oplus B,C) = (\dim A + \dim B)\dim C\) carries a
\(B\)-independent term \((\dim A)(\dim C)\) that no single object
\([A,C]_\oplus\) can absorb across all \(B\), so no right adjoint
exists. For \(\otimes\),
\(\dim\mathrm{Hom}(A\otimes B,C) = (\dim B)\cdot(\dim A\dim C)\) is
proportional to \(\dim B\) and is absorbed by \(A\multimap C\). In
short: \(\otimes\) is \emph{multiplicative} (dimensions multiply,
currying exists, closed but no diagonal); \(\oplus\) is \emph{additive}
(dimensions add, diagonal exists, but not closed). Diagonal and eval
live on different monoidal products and do not coincide.

\hypertarget{extending-the-transformer-to-satisfy-the-lnl-adjunction-and-its-cost}{%
\subsection{Extending the Transformer to satisfy the LNL adjunction, and its cost}\label{extending-the-transformer-to-satisfy-the-lnl-adjunction-and-its-cost}}

Benton's LNL model is a monoidal adjunction \(F\dashv G\) between a CCC \(\mathcal{C}\) (cartesian, \(\times,\Rightarrow\)) and an SMCC \(\mathcal{L}\) (linear, \(\otimes,\multimap\)), with \(!=F\circ G\).
The Transformer's parts map onto the two worlds --- \(\mathcal{L}=\) the value-path SMCC plus the softmax Markov category; \(\mathcal{C}=\) the ReLU MLP (a PL/cartesian fragment) --- but the \emph{adjunction itself
is absent}: the MLP is a self-morphism within one category, not a functor pair \(F:\mathcal{C}\rightleftarrows\mathcal{L}:G\) between the
two worlds. So \textbf{the Transformer does not satisfy the LNL adjunction}; this is the LNL restatement of ``\(!\) is absent.''

Extending it to satisfy LNL means implementing \(F\dashv G\): a read-out
\(G:\mathcal{L}\to\mathcal{C}\) (linear vectors to copyable discrete/PL representations) and an embedding \(F:\mathcal{C}\to\mathcal{L}\), whose
composite \(!=F\circ G\) is an explicit duplication mechanism --- write an intermediate value out to a copy-free store, reuse it freely, read it
back. With \(!\) in place, linear logic embeds intuitionistic logic (\(A\to B = {!A}\multimap B\)), so intermediate results become unboundedly reusable and the sequential state-tracking that the plain model cannot express (§9.4) becomes reachable in principle.

The extension carries three costs, each sharpened by our results.

\begin{itemize}
\item
  \textbf{Loss of parallelism.} Unbounded reuse via \(F\circ G\) is
  sequential computation, incompatible with the constant-depth parallel
  form that puts the model in \(\mathsf{TC}^0\) (§9.4,
  Merrill--Sabharwal). Satisfying LNL pushes the system out of
  \(\mathsf{TC}^0\) toward \(\mathsf{NC}^1\) --- i.e.~toward the RNN,
  not a faster Transformer.
\item
  \textbf{Inconsistency with the depth-as-grade experiment.} A
  depth-\(L\) model can unfold \(F\circ G\) at most \(L\) times, giving
  only a graded \(!_L\), not a full \(!\). But §9.4 and the \((r,L)\) /
  three-architecture sweeps \emph{falsify} the reading that depth
  supplies a per-reuse budget: the state-tracking boundary \(n^*\) is
  flat in depth (and in heads), a fixed structural constant, not
  \(n^*\sim c\cdot 2^L\). So the naïve ``depth \(=\) number of LNL
  unfoldings'' extension lacks empirical support; adding \(!_L\)
  internally is not what the data show the architecture doing.
\item
  \textbf{Lossy read-out.} \(G\) (linear \(\to\) cartesian) must discretise, discarding the continuous/probabilistic content (routes 3--4); the extension partially kills the very Markov structure that §6
  identifies.
\end{itemize}

The upshot: satisfying LNL trades the Transformer's defining parallelism
for sequential reuse, and the most faithful realisation is not an
internally-modified Transformer but a \textbf{hybrid} --- a
linear/parallel controller (\(\mathcal{L}\)) coupled to a
cartesian/sequential external store (\(\mathcal{C}\)) --- which is
exactly the chain-of-thought and tool-use regime of §10.5.

\hypertarget{finite-dimensionality-does-not-refute-the-attention-matrix-is-a-power-object-element}{%
\subsection{Finite-dimensionality does not refute ``the attention matrix is a power-object element''}\label{finite-dimensionality-does-not-refute-the-attention-matrix-is-a-power-object-element}}

Critiques of Turing-completeness claims for Transformers rest on finite
dimensionality and finite precision: fixed-size vectors cannot encode an
unbounded tape, so a fixed Transformer is not a universal machine. These
critiques are correct, but they do not touch the present paper, because
our claims are about a \emph{different property}. The reading
``attention matrix \(A\) realises an element of the internal hom, and
apply \(=\) eval'' (§3, §7) is defined entirely in a
\textbf{finite-dimensional SMCC}: \(\mathbf{FinVect}\) is symmetric
monoidal closed, with internal hom \(A\multimap B\) of dimension
\((\dim A)(\dim B)\) and evaluation
\(\mathrm{ev}:(A\multimap B)\otimes A\to B\) a finite-dimensional linear
map. The linear-\(\lambda\) reading is a \emph{type system}, not a
demand for unbounded computational resource; indeed this paper's thesis
is the \emph{absence} of \(!\) --- bounded, resource-linear computation.
Finite-dimensionality is therefore not an obstruction but the natural
home of the eval structure (infinite-dimensional spaces make the
internal hom worse-behaved, not better). Our account already sits on the
\emph{finite/bounded} side: §5's no-\(!\), §9.4's \(\mathsf{TC}^0\)
ceiling. The Turing-completeness critique and this paper are on the same
side of the ledger, not in conflict. (The one genuine conditionality of
``\(A\) = power-object element'' is not dimensionality but the
representability of §4b --- the strength of \(D\) --- left as an
obligation in the Lean development; finite dimensionality in fact makes
that strength easier to exhibit, not harder.)

\hypertarget{infinite-external-memory-universality-and-its-categorical-meaning}{%
\subsection{Infinite external memory, universality, and its categorical meaning}\label{infinite-external-memory-universality-and-its-categorical-meaning}}

If the scratchpad is unbounded, the \emph{system} can become
Turing-universal --- but the subject of that predicate is not the
Transformer. A finite controller plus an unbounded tape is already a
Turing machine; a Transformer (a finite-dimensional, \(\mathsf{TC}^0\)
finite control) plus an unbounded external store is universal for the
same classical reason, and, as the critiques note, this establishes
universality of the \emph{system}, not of the network --- just as a
finite-state machine with a tape is Turing-complete without the FSM being so. In our terms this is the limit of the LNL story of §10.3: the
unbounded store is the cartesian world \(\mathcal{C}\) (copy-free tape),
and read/write is \(G/F\); the controller (\(\mathcal{L}\), parallel, \(\mathsf{TC}^0\)) acquires \(!\) --- unbounded reuse --- \emph{at the level of the system} by externalising it, rather than internalising it
and losing parallelism. This is a coherent design: keep the network
parallel and cheap, delegate sequentiality and unbounded reuse to the
store, and couple them by an LNL-like read/write. It is exactly the
tool-use / code-execution / retrieval setting in which large models
actually operate --- with the caveat (which we adopt) that the coupling
is engineered, lossy, and does not verify the strict adjunction laws, so
it is a system-level \emph{approximation} of \(!\)-externalisation, not an internal LNL model.

\hypertarget{relation-to-merrillsabharwal-2024-chain-of-thought}{%
\subsection{Relation to Merrill--Sabharwal 2024 (chain-of-thought)}\label{relation-to-merrillsabharwal-2024-chain-of-thought}}

Merrill and Sabharwal \cite{merrill2024expressive} show that chain-of-thought raises the
expressive power of Transformers above \(\mathsf{TC}^0\), scaling with the number of generated intermediate steps. Our framework offers a
categorical reading of \emph{why}: CoT is the sequential unfolding of the LNL read-out \(G:\mathcal{L}\to\mathcal{C}\) --- writing
intermediate values out as tokens is exactly moving into the copy-free
cartesian world, from which they can be reused, i.e.~an \textbf{externalised \(!\)}. This is the same relation our §9.4 has to Merrill et al.'s \(\mathsf{TC}^0\) classification\cite{merrill2023parallelism}: they characterise the boundary from the outside (circuit complexity), we supply the inside (which categorical structure crosses it).

Two points must be stated precisely to avoid over-claiming. First, this is an \emph{interpretive} strengthening: we do not re-derive the CoT
theorem categorically, and the CoT-as-\(G\)-unfolding correspondence inherits the caveats above (not a strict adjunction, lossy, external).
Second --- and this corrects a tempting misreading --- our falsification of depth-as-grade (§10.3) does \textbf{not} assert a limitation of CoT.
Depth is an \emph{internal} axis; CoT is an \emph{external} one. That internal depth fails to supply a reuse budget is not evidence against
the external mechanism; on the contrary, ``\(!\) cannot be grown internally (by depth)'' is precisely what makes the
\emph{externalisation} of \(!\) (CoT, tools) necessary. The correct statement is a division of labour: the exponential modality is
unavailable from the Transformer's internal structure (depth, width, heads all leave \(n^*\) flat, §9.4) and is instead supplied, at the
system level, by external sequential generation --- which is the categorical content of chain-of-thought's expressivity gain.

\hypertarget{learnability-of-transformers}{%
\subsection{Learnability of transformers}\label{learnability-of-transformers}}

The success of transformers is not only higher order function programmability and in-context learning, but learnablity and avoiding
local minimum, overfitting are also significant properties and affect to large application area industries. For example, Edge of chaos hypothesis
states highset learning speed is achieved when learning rate is on critical point\cite{EoC}. In other studies, attentions as a component of transformers tends to cluster in reccuerent
structure\cite{KuramotoTransformer}. On the other hand MLP suffers from
chaotic separation of phase spate\cite{MagicNumber7} which cause poor
classification resulst. As combinations of attention and MLP,
transformers can be adjust learning dynamics properly to reach low loss
function solution speedy. In this case changing the ration between
attentions and MLPs and measure prediction performance of learned
parameters is simple experiment to detect the function of edge of chaos flow\cite{EoC}.

In this paper, we only show formulation and explanation of inference and generation of fuction Transformers and ability higher logical property.
To extend categoric theoretical view to the learnability of transformer, explaining this dynamical systems point of views are required. Because learning process cannot treated as natural transformation. Actually, 2-category which has morphism of morphism as an objcet, is nesesssary to explain training process.

On the other hand, ICL can ben treated as linear \(\lambda\)-calculas based on linear category. Where the difference comes from is interesting theme.

Using Cartesian reverse derivative category(CRDC) which has is an answer\cite{CRDC}. Discriminaiton 2 kinds of Jacobian \(R[f]\)(partial
derivativs of weight parameter) and \(D_A[f]\)(derivatives of layer input vector), are different variables but they are connected with chain
rune of differnetation. Lyapunov spectrum ,eigen values of these Jacobians rule dynamics of neural networks. Differencial structure,
spectrum and topology are additional structure of CRDC and can be described by using vocablaries of basic category theory such as functor or natural transformation.

Whether the extra structure reduces to categories/functors/natural transformations in the same sense that 2-categories are \(\mathbf{Cat}\)-enriched and \(\infty\)-categories are \(\mathbf{sSet}\)-enriched:

\begin{longtable}[]{@{}llll@{}}
\toprule
\begin{minipage}[b]{0.22\columnwidth}\raggedright
Added structure\strut
\end{minipage} & \begin{minipage}[b]{0.22\columnwidth}\raggedright
Basic-vocabulary formulation\strut
\end{minipage} & \begin{minipage}[b]{0.22\columnwidth}\raggedright
Ground required externally\strut
\end{minipage} & \begin{minipage}[b]{0.22\columnwidth}\raggedright
Reducibility\strut
\end{minipage}\tabularnewline
\midrule
\endhead
\begin{minipage}[t]{0.22\columnwidth}\raggedright
Differential\strut
\end{minipage} & \begin{minipage}[t]{0.22\columnwidth}\raggedright
Tangent category: functor \(T\) + natural transformations
\((p,0,+,c,l)\) + limit preservation\strut
\end{minipage} & \begin{minipage}[t]{0.22\columnwidth}\raggedright
(none essential)\strut
\end{minipage} & \begin{minipage}[t]{0.22\columnwidth}\raggedright
Fully internal\strut
\end{minipage}\tabularnewline
\begin{minipage}[t]{0.22\columnwidth}\raggedright
Spectral\strut
\end{minipage} & \begin{minipage}[t]{0.22\columnwidth}\raggedright
Dagger category + biproducts (contravariant \(\dagger\), natural
isos)\strut
\end{minipage} & \begin{minipage}[t]{0.22\columnwidth}\raggedright
Scalar object \(\Bbbk=\mathbb{C}\), algebraically closed and
complete\strut
\end{minipage} & \begin{minipage}[t]{0.22\columnwidth}\raggedright
Structure reducible; eigenvalue \emph{content} is ground\strut
\end{minipage}\tabularnewline
\begin{minipage}[t]{0.22\columnwidth}\raggedright
Topological equivalence\strut
\end{minipage} & \begin{minipage}[t]{0.22\columnwidth}\raggedright
\(\mathbf{Top}\)-enrichment + preservation of the \(\mathbb{R}\)-action
(flow)\strut
\end{minipage} & \begin{minipage}[t]{0.22\columnwidth}\raggedright
Base \(\mathbf{Top}\) (or condensed) and a time object\strut
\end{minipage} & \begin{minipage}[t]{0.22\columnwidth}\raggedright
Only via enrichment; not internalized\strut
\end{minipage}\tabularnewline
\bottomrule
\end{longtable}

So all three are expressible with categories, functors and natural
transformations but, exactly as for 2- and \(\infty\)-categories,
spectral theory and topological equivalence require \emph{choosing an
enrichment base} (\(\mathbb{C}\)-linear complete dagger;
\(\mathbf{Top}\)). Only the differential layer is purely internal
(tangent categories). The minimal categorical setting for the
bifurcation program is therefore: tangent category (differential,
internal) + \(\mathbb{C}\)-linear complete dagger-biproduct enrichment
(spectral, chosen ground) + \(\mathbf{Top}\)-enrichment or flow
\(\mathbb{R}\)-action (topological, external). The most tractable route
is to phrase grokking's saddle-to-saddle as a loss of invertibility /
imaginary-axis crossing of the state Jacobian \(D_A[f]\) in the dagger
subcategory, which uses only the first two layers; topological conjugacy
failure then follows via Hartman--Grobman.

In general, considering network architectures in points of view of
dynamical systems and category theory is useful for the performance and
its limit. Especially CRDC relates to and describe lyapunov spectrum ,
bifurcation, learning dynamics like grokking is important question.

\hypertarget{conclusion}{%
\section{Conclusion}\label{conclusion}}

A Transformer layer is faithfully described not by one symmetric
monoidal closed category but by a composite: a Markov category (softmax)
bridged by a \(D\)-algebra expectation to an SMCC (value path, where
evaluation and a linear-λ reading live), followed by a non-linear
realization map (MLP) and threaded by an additive residual copy. Copy
exists in two legitimate but distinct forms --- additive (residual,
depth axis) and Markov (softmax, position axis) --- neither of which is
the exponential modality, so the layer is resource-linear at the
evaluation site. Controlled experiments localize the copy structure in
softmax and the applied value in the MLP; the Function-Vector
experiments that would confirm this at language-model scale are
specified and left as placeholders.

On expressivity we defer to the topos analysis of Villani and McBurney \cite{Topos}: for ReLU networks the PL base is apt, and their choose-eval
decomposition independently reaches the eval--apply reading at the architecture level. What we add is a different axis. 
In the linear/PL-input regime, the value path --- residual copy included --- is expressible in \(!\)-free intuitionistic linear λ-calculus, and the hypothesis that the model performs \emph{only} such resource-linear computation cannot be excluded; the experiments of §8 support rather than merely fail to refute it.

The same resource lens reads the \emph{limits} of the architecture
(§9.4): on permutation-composition state tracking, a shared-curriculum
three-architecture sweep reproduces the \(\mathsf{TC}^0\) classification
of Merrill and collaborators as three capacity-independent plateaus (RNN
high, Transformer intermediate, diagonal SSM at the floor), and the
plateau heights track the copy structure of each forward category ---
sequential exponential (RNN), bounded parallel Markov coupling
(Transformer), none (SSM). The Transformer's intermediate foothold
requires both the softmax Markov kernel and the cartesian MLP
nonlinearity, neither alone.

\textbf{Open problems.} (i) Whether the non-natural Markov copy of
softmax effectively supplies duplication of the strength of \(!\) ---
testable by asking whether reuse of a single FV across multiple queries
depends on attention temperature/sharpness. (ii) Which component
realizes \(\Phi\) in a real LM (MLP vs the \(QK^\top\) currying), via
the FV tests of §8.2. (iii) Whether the PL/topos higher-order account
and the \(!\)-free linear account can be reconciled as a single
stratified model: a higher-order \emph{architecture-selection} level over a resource-linear \emph{value-application} level.

\hypertarget{aknowledgement}{%
\subsection{Aknowledgement}\label{aknowledgement}}

We thank very useful discussion with Dr. Kunihiko Kaneko and Dr. Kai Nakaishi.

\hypertarget{references}{%
\section{References}\label{references}}
\bibliographystyle{plain}
\bibliography{fmap_cat}

\emph{Software:} \texttt{ericwtodd/function\_vectors}; the synthetic-probe script \texttt{eval\_apply\_probes.py} and the FV-native script \texttt{eval\_apply\_fv.py} accompanying this note.

\hypertarget{appendix}{%
\section{Appendix}\label{appendix}}

\hypertarget{appendix-a-formal-proofs-of-theorems-in-lean}{%
\subsection{Appendix A formal proofs of theorems in lean}\label{appendix-a-formal-proofs-of-theorems-in-lean}}

Common sector

\begin{lstlisting}
import Mathlib.CategoryTheory.Monoidal.Closed.FunctorCategory.Basic
import Mathlib.CategoryTheory.Monoidal.Closed.Basic
import Mathlib.CategoryTheory.Monoidal.Braided.Basic
import Mathlib.CategoryTheory.Limits.Shapes.BinaryBiproducts
import Mathlib.CategoryTheory.Preadditive.Biproducts
import Mathlib.CategoryTheory.Limits.Shapes.IsTerminal
import Mathlib.CategoryTheory.Monad.Kleisli
import Mathlib.CategoryTheory.Monad.Basic
import Mathlib.CategoryTheory.Monad.Algebra
import Mathlib.CategoryTheory.Adjunction.Basic

open CategoryTheory CategoryTheory.MonoidalCategory CategoryTheory.Limits

noncomputable section
namespace TransformerCat

universe v u
variable {C : Type u} [Category.{v} C]
\end{lstlisting}

``noncomputable section'' means following theorems do not caluculated
specific values as functions.

\begin{itemize}
%\tightlist
\item
  Theorem 1 The softmax routing is a Kleisli morphism of a monad D, and   whole layer of transformer is Kleisli morphism of composit monado M =  $T \cdot D$.
\end{itemize}

\begin{lstlisting}
section KleisliLayer
variable (T D : Monad C)
variable {A B : C}

/-- In `Kleisli D`, a morphism A → B is by definition a base morphism
    A → D.obj B. The softmax/Markov routing `attn : A → D.obj B` (D = 
    the distribution monad) is therefore literally a Kleisli morphism of D. -/
example (attn : A → (D : C => C).obj B) :
    @Quiver.Hom (Kleisli D) _ (A : Kleisli D) (B : Kleisli D) := attn

/-- Data representing a distributive law together with the composite monad
that its omitted Beck coherence equations are intended to induce. -/
/- The original blueprint used `True` in place of Beck's four coherence
axioms. Those placeholders cannot justify construction of a composite monad.
Until those equations are formalized, the honest interface must include the
resulting monad as data. -/
structure DistribLaw (T D : Monad C) where
  law : (D : C => C) >>> (T : C => C) -> (T : C => C) >>> (D : C => C)
  composite : Monad C
  composite_toFunctor : (composite : C => C) = (T : C => C) >>> (D : C => C)

/-- The composite monad supplied by `DistribLaw`. -/
def composeMonad (l : DistribLaw T D) : Monad C := l.composite

/-- **The whole layer as a single Kleisli morphism of the composite monad.**
    Given the composite monad M = `composeMonad`, a layer
    `layer : A → M.obj B` is exactly a morphism `A → B` in `Kleisli M`.
    Thus the two categories (Markov `Kl(D)` and the value/SMCC part carried by T)
    are unified as morphisms of the single Kleisli category `Kleisli M`. -/
example (l : DistribLaw T D)
    (layer : A → ((composeMonad T D l) : C => C).obj B) :
    @Quiver.Hom (Kleisli (composeMonad T D l)) _
      (A : Kleisli (composeMonad T D l)) (B : Kleisli (composeMonad T D l)) :=
  layer

end KleisliLayer
\end{lstlisting}

\begin{itemize}
%\tightlist
\item
  Theorem 2 eval-apply is unit/counit of adjoint, the type of
  \(\lambda\)c.\(\lambda\)x.(Φ(Ec))x is linear \(\lambda\) term.
\end{itemize}

\begin{lstlisting}
section EvalApply
variable [MonoidalCategory C] [MonoidalClosed C]
variable {Ctx P X Y : C}

/- Extraction E : Ctx → P_FV and realization Φ : P_FV → (X ⊸ Y).
   Here `(ihom X).obj Y` is the internal hom X ⊸ Y. -/
variable (E : Ctx → P) (Φ : P → (ihom X).obj Y)

/-- The reified, realized function  f_t = Φ ∘ E : Ctx → (X ⊸ Y).
    This is the (linear) $\lambda$-abstraction / "choose" morphism. -/
def curriedFn : Ctx → (ihom X).obj Y := E ≫ Φ

/-- Application: uncurrying the reified function, X ⊗ Ctx → Y. -/
def applyMor : X ⊗ Ctx → Y := MonoidalClosed.uncurry (E ≫ Φ)

/-- **eval-apply.**  Applying the reified function equals "build it (X ◁ (E ≫ Φ)),
    then evaluate", where the evaluation `ihom.ev X` is exactly the counit of the
    adjunction `(tensorLeft X) ⊣ (ihom X)`.  This is the categorical content of
    `(Φ (E c)) x = eval (Φ (E c), x)`. -/
theorem apply_eq_build_then_ev :
    applyMor E Φ = (X ◁ (E ≫ Φ)) ≫ (ihom.ev X).app Y := by
  unfold applyMor
  rw [MonoidalClosed.uncurry_eq]

/-- **β-conversion** = the counit triangle: uncurry (curry g) = g. -/
theorem beta (g : X ⊗ Ctx → Y) :
    MonoidalClosed.uncurry (MonoidalClosed.curry g) = g :=
  MonoidalClosed.uncurry_curry g

/-- **η-conversion** = the unit triangle: curry (uncurry h) = h. -/
theorem eta (h : Ctx → (ihom X).obj Y) :
    MonoidalClosed.curry (MonoidalClosed.uncurry h) = h :=
  MonoidalClosed.curry_uncurry h

/-
  The term  $\lambda$c. $\lambda$x. (Φ (E c)) x  of linear $\lambda$-calculus, of type
  Ctx ⊸ (X ⊸ Y), DENOTES `curriedFn E Φ : Ctx → (ihom X).obj Y`, and its
  applied form denotes `applyMor E Φ`. Theorems `apply_eq_build_then_ev`, `beta`,
  `eta` are the semantic (SMCC) counterparts of the term's typing plus β/η.

  A *syntactic* soundness theorem — "this linear-$\lambda$ term is well-typed with each
  variable used exactly once, and its denotation is `applyMor`" — requires a
  formalized linear type system (contexts as multisets, ⊸-intro/elim, a
  no-contraction/no-weakening discipline) that Mathlib does NOT provide. That is
  a separate development; here we formalize the denotation only.
-/

/-- Naturality bookkeeping: uncurrying commutes with precomposition by E,
    i.e. the "choose then apply" pipeline composes as expected. -/
theorem apply_factor :
    applyMor E Φ = (X ◁ E) ≫ MonoidalClosed.uncurry Φ := by
  unfold applyMor
  rw [MonoidalClosed.uncurry_natural_left]

end EvalApply
\end{lstlisting}

\begin{itemize}
\item
  Theorem 3 residual connection is written by (co)diagonal of biproduct, this is not !.
\end{itemize}

\begin{lstlisting}
section Residual
variable [Preadditive C] [HasBinaryBiproducts C]
variable {A : C}

/-- The additive diagonal Δ_⊕ : A → A ⊞ A (fan-out along the depth axis). -/
def diagAdd (A : C) : A → A ⊞ A := biprod.lift (𝟙 A) (𝟙 A)

/-- The additive codiagonal ∇_⊕ : A ⊞ A → A (write-back). -/
def codiagAdd (A : C) : A ⊞ A → A := biprod.desc (𝟙 A) (𝟙 A)

/-- **Residual as additive copy (clean form).**
    `biprod.lift (𝟙) f ≫ biprod.desc (𝟙) (𝟙) = 𝟙 + f`.
    The two branches Δ_⊕ produces are summed back by ∇_⊕ into a SINGLE
    resource `𝟙 + f`; this is why the residual copies additively but supplies
    no independent second consumption. -/
theorem residual_eq (f : A → A) :
    biprod.lift (𝟙 A) f ≫ biprod.desc (𝟙 A) (𝟙 A) = 𝟙 A + f := by
  simp [biprod.lift_desc]

/-
**Residual as Δ_⊕ ≫ (id ⊞ f) ≫ ∇_⊕.**
    Same statement, written through the diagonal / map / codiagonal, matching
    the string-diagram reading.
-/
theorem residual_eq_diag (f : A → A) :
    diagAdd A ≫ biprod.map (𝟙 A) f ≫ codiagAdd A = 𝟙 A + f := by
  simp +decide [ diagAdd, codiagAdd, ← Category.assoc ];
  grind +suggestions

/-
The original proposed theorem `no_tensor_diagonal_of_noncartesian` was
incorrect: in a preadditive monoidal category the family of zero maps is always
such a natural diagonal.  A counit, including its normalization and naturality,
is needed to derive the advertised obstruction.

A natural, normalized family of discarding maps would make the tensor unit
terminal. Hence such a family cannot exist when the tensor unit is not
terminal. This is the part of the obstruction to a uniform copying/discarding
comonoid structure that follows directly from non-cartesianness.
-/
omit [Preadditive C] [HasBinaryBiproducts C] in
theorem no_natural_discard_of_nonterminal_unit
    [MonoidalCategory C]
    (hNonterminal : IsEmpty (Limits.IsTerminal (𝟙_ C))) :
    ¬ ∃ ε : (A : C) → (A → 𝟙_ C),
        ε (𝟙_ C) = 𝟙 (𝟙_ C) ∧
        (∀ {A B : C} (g : A → B), g ≫ ε B = ε A) := by
  rintro ⟨ε, hunit, natural⟩
  apply hNonterminal.false
  refine Limits.IsTerminal.ofUniqueHom ε ?_
  intro X m
  simpa [hunit] using natural m

end Residual
\end{lstlisting}

\begin{itemize}
%\tightlist
\item
  Theorem 4 The two roles of the attention matrix, \(Kl(D)(pos,pos)\)
  and \(Hom(V\multimap V)\) connected by a functor.
\end{itemize}

\begin{lstlisting}
section RepresentationFunctor
variable (D : Monad C)

/-- **The representation functor `F_V` (Type-valued), FULLY PROVED functorial.**
    A probability kernel `A` is sent to the value-mixing operator on `pos → V`.
    `map_id` uses the unit of the monad and of the algebra; `map_comp` uses the
    multiplication, its naturality, and the algebra's associativity. No strength, no `sorry`. -/
def valuePresheaf (Valg : D.Algebra) : (Kleisli D)ᵒᵖ => Type v where
  obj X := X.unop → Valg.A
  map {X Y} A := fun val => A.unop ≫ (D : C => C).map val ≫ Valg.a
  map_id X := by
    funext val
    simp only [unop_id]
    -- Kleisli identity is the monad unit η; then η-naturality + algebra unit.
    show D.η.app X.unop ≫ (D : C => C).map val ≫ Valg.a = val
    rw [← Category.assoc, ← D.η.naturality val, Category.assoc, Valg.unit,
        Category.comp_id]
  map_comp {X Y Z} A B := by
    funext val
    -- opposite comp unops to reversed Kleisli comp
    --   (A ≫ B).unop = B.unop ≫_Kl A.unop = B.unop ≫ D.map A.unop ≫ μ ;
    -- expand D.map of the composite on the right, then μ-naturality + algebra assoc.
    show (B.unop ≫ (D : C => C).map A.unop ≫ D.μ.app X.unop)
            ≫ (D : C => C).map val ≫ Valg.a
        = B.unop ≫ (D : C => C).map (A.unop ≫ (D : C => C).map val ≫ Valg.a) ≫ Valg.a
    rw [Functor.map_comp, Functor.map_comp]
    simp only [Category.assoc]
    rw [D.μ.naturality_assoc, Valg.assoc]

/-- **Functoriality made explicit: the kernel action respects Kleisli identity.**
    `A = η` (the deterministic "stay put" kernel) acts as the identity operator. -/
theorem valuePresheaf_map_id (Valg : D.Algebra) (X : (Kleisli D)ᵒᵖ) :
    (valuePresheaf D Valg).map (𝟙 X) = id :=
  (valuePresheaf D Valg).map_id X

/-- **and respects Kleisli composition (Chapman-Kolmogorov ↦ operator comp).** -/
theorem valuePresheaf_map_comp (Valg : D.Algebra) {X Y Z : (Kleisli D)ᵒᵖ}
    (A : X → Y) (B : Y → Z) :
    (valuePresheaf D Valg).map (A ≫ B)
      = (valuePresheaf D Valg).map A ≫ (valuePresheaf D Valg).map B :=
  (valuePresheaf D Valg).map_comp A B

end RepresentationFunctor
\end{lstlisting}


\hypertarget{appendix-b-the-relation-between-distribution-moand-d-kleisli-categgory-mathrmkld-and-d-algebra-category-mathrmemd}{%
\subsection{\texorpdfstring{Appendix B: The relation between distribution moand D, Kleisli categgory \(\mathrm{Kl}(D)\) and D-algebra category
\(\mathrm{EM}(D)\)}{Appendix B: The relation between distribution moand D, Kleisli categgory \textbackslash mathrm\{Kl\}(D) and D-algebra category \textbackslash mathrm\{EM\}(D)}}\label{appendix-b-the-relation-between-distribution-moand-d-kleisli-categgory-mathrmkld-and-d-algebra-category-mathrmemd}}

\hypertarget{the-distribution-monad-d}{%
\subsubsection{\texorpdfstring{The Distribution Monad \(D\)}{The Distribution Monad D}}\label{the-distribution-monad-d}}

Let the base category be \(\mathbf{Set}\), and define the finite-support version (the continuous version is discussed later):
\[D(X)=\Big\{\,p:X\to[0,1]\ \Big|\ \mathrm{supp}(p)\text{ finite},\ \textstyle\sum_{x}p(x)=1\,\Big\}\]

This is the set of finitely-supported probability distributions on \(X\). The three-part package that makes it a monad:

\begin{itemize}
%\tightlist
\item
  \textbf{Unit} \(\eta_X:X\to D(X)\), \(x\mapsto \delta_x\) (the Dirac measure, point mass). ``Regard a definite value as a distribution.''
\item
  \textbf{Multiplication} \(\mu_X:D(D(X))\to D(X)\), collapsing a distribution of distributions by averaging:
  \(\mu(P)(x)=\sum_{q}P(q)\,q(x)\). This is exactly the \textbf{law of total probability}.
\item
  \textbf{Functorial action} \(D(f):D(X)\to D(Y)\) (for \(f:X\to Y\)) is the \textbf{pushforward}: \(D(f)(p)(y)=\sum_{x:f(x)=y}p(x)\).
\end{itemize}

The monad laws (left/right unit laws and associativity) coincide with the \textbf{basic identities of probability}: the marginalization of
Dirac measures and the associativity of mixing. Here \(\eta\) is ``making definite'' and \(\mu\) is ``flattening a mixture.''

\textbf{On the base category}: \(D\) is \textbf{commutative} (the sampling order of two independent distributions does not change the
result = Fubini) and \textbf{affine} (\(D(1)\cong 1\); there is exactly one distribution on a one-point set). These two properties are what
later make \(\mathrm{Kl}(D)\) a Markov category.

\textbf{Variants}: the \textbf{Giry monad} \(G(X)=\{\)probability measures on \(X\}\) over measurable spaces \(\mathbf{Meas}\) (the
continuous version --- this is the right one for the real-valued logits of attention); the subdistribution monad for subprobabilities
(\(\sum\le 1\)); and so on. Since softmax has finite support, \(D\) suffices.

\hypertarget{the-kleisli-category-mathrmkld-the-category-of-stochastic-kernels}{%
\subsubsection{\texorpdfstring{The Kleisli Category \(\mathrm{Kl}(D)\)
--- the Category of Stochastic Kernels}{The Kleisli Category \textbackslash mathrm\{Kl\}(D) --- the Category of Stochastic Kernels}}\label{the-kleisli-category-mathrmkld-the-category-of-stochastic-kernels}}

\(\mathrm{Kl}(D)\) is the category that puts ``stochastic morphisms'' center stage. - \textbf{Objects}: the same as the base category (sets).
- \textbf{Morphisms} \(X\to Y\): functions \(X\to D(Y)\),
i.e.~assignments of a distribution on \(Y\) to each \(x\) --- \textbf{Markov kernels (stochastic kernels)}. For a finite set, a \textbf{stochastic matrix} \(k(x)(y)=P(y\mid x)\ge0\) whose columns (or rows) sum to 1. - \textbf{Identity morphism}: \(\eta_X\),
\(x\mapsto\delta_x\) (deterministically pass through unchanged). -\textbf{Composition} (Kleisli composition = \textbf{Chapman--Kolmogorov}):
\[(k'\circ k)(x)(z)=\sum_{y}k(x)(y)\,k'(y)(z)\] the product of stochastic matrices. ``Marginalize over the intermediate \(y\) the transition from \(x\) to \(y\) and from \(y\) to \(z\).'' - \textbf{Symmetric monoidal structure} (from the commutativity of \(D\)):
\(\otimes=\) the Cartesian product of sets, and the tensor of kernels = the independent joint distribution. - \textbf{Markov category
structure}: copy \(X\to X\otimes X\), \(x\mapsto\delta_{(x,x)}\) (deterministic duplication) and delete \(X\to 1\). Because \(D\) is
affine, delete is unique (semicartesian); because it is commutative, it is symmetric monoidal. \textbf{This copy is non-natural} --- for a
genuinely stochastic kernel \(k\), ``copy then \(k\)'' (a perfectly correlated pair \((y,y)\)) does not agree with ``\(k\) then copy''
(independent resampling). It is a copy that generates correlation.

Restricting \(\mathrm{Kl}(D)\) to finite sets gives \textbf{FinStoch} (the category of stochastic matrices). \textbf{softmax\((QK^\top)\) is
precisely a morphism of this category} --- a stochastic kernel from query positions to key positions, a stochastic matrix \(A\).

\hypertarget{d-algebras-the-eilenbergmoore-category-mathrmemd}{%
\subsubsection{\texorpdfstring{\(D\)-Algebras (the Eilenberg--Moore Category \(\mathrm{EM}(D)\))}{D-Algebras (the Eilenberg--Moore Category \textbackslash mathrm\{EM\}(D))}}\label{d-algebras-the-eilenbergmoore-category-mathrmemd}}

Whereas Kleisli made ``morphisms'' the protagonist, \(D\)-algebras make ``\textbf{a structure that consumes a distribution and returns a value}'' the protagonist.

A \textbf{\(D\)-algebra} is a pair \((X,\alpha)\) with
\(\alpha:D(X)\to X\) satisfying two coherence laws:
the point mass $\delta_x$ evaluates to x and a distribution of distributions (flatten then evaluate = evaluate each then evaluate)

\[\alpha\circ\eta_X=\mathrm{id}_X\qquad \]
\[\alpha\circ\mu_X=\alpha\circ D(\alpha)\qquad \]

\textbf{Meaning}: \(\alpha\) maps ``a distribution on \(X\) to a single element,'' i.e.~it performs \textbf{taking a barycenter / expectation / convex combination}. In other words, a \(D\)-algebra = a set on which convex combinations can be taken.

\textbf{What concretely is a \(D\)-algebra}: the Eilenberg--Moore category of the finite-distribution monad is equivalent to the category of \textbf{convex spaces (abstract convex spaces)}\cite{Markov}. The decisive example for us --- \textbf{any real vector space \(V\) is a \(D\)-algebra}, with structure map

\[\alpha_V:D(V)\to V,\qquad \alpha_V(p)=\sum_{v}p(v)\,v=\mathbb E_{p}[v]\quad(\textbf{expectation}).\]

More generally, a convex subset of a vector space is a \(D\)-algebra.
The morphisms of algebras are affine maps (commuting with convex combinations = commuting with \(\alpha\)).

\textbf{Value mixing is exactly this action}: Because after softmax
produces \(A(x)\in D(\text{positions})\), the value mixing \(A\!\cdot\!V=\sum_y A(x)(y)\,v_y\) is the \(D\)-algebra structure
\(\alpha_V:D(V)\to V\) of the value space \(V\) applied to the distribution = the expectation.

\[\text{attention head}=\underbrace{\big[\mathrm{Ctx}\xrightarrow{\ \text{softmax}\ }D(\text{pos})\big]}_{\text{morphism of }\mathrm{Kl}(D)}\ \text{composed with}\ \underbrace{\big[D(V)\xrightarrow{\ \alpha_V=\mathbb E\ }V\big]}_{\text{action of a }D\text{-algebra}}.\]

Consistency with Kolmogorov theory: a random variable = a measurable function on the sample space \(\Omega\to\mathbb R\), and the expectation \(\mathbb E\) is the \(D\)-algebra (Giry-algebra) structure of \(\mathbb R\) (or \(V\)). \textbf{softmax routes and expectation mixes}
fits into the single phrase ``kernel $\cdot$ algebra action.''

\hypertarget{the-relationship-between-mathrmkld-and-mathrmemd-adjunctions-and-the-comparison-functor}{%
\subsubsection{\texorpdfstring{The Relationship Between \(\mathrm{Kl}(D)\) and \(\mathrm{EM}(D)\) --- Adjunctions and the Comparison Functor}{The Relationship Between \textbackslash mathrm\{Kl\}(D) and \textbackslash mathrm\{EM\}(D) --- Adjunctions and the Comparison Functor}}\label{the-relationship-between-mathrmkld-and-mathrmemd-adjunctions-and-the-comparison-functor}}

The two are not unrelated; they are two resolutions of the same monad \(D\). Every monad arises from an adjunction \(F\dashv U\), and \(D\) has two canonical ones.

\begin{itemize}
%\tightlist
\item
  \textbf{Kleisli resolution}: \(F_K:\mathbf{Set}\to\mathrm{Kl}(D)\)
  (free) \(\dashv U_K\). \(\mathrm{Kl}(D)\) is the category of
  \textbf{free \(D\)-algebras} only.
\item
  \textbf{Eilenberg--Moore resolution}:
  \(F_{EM}:\mathbf{Set}\to\mathrm{EM}(D)\), \(X\mapsto(DX,\mu_X)\) (free
  algebra) \(\dashv U_{EM}\) (forgetful, \((X,\alpha)\mapsto X\)).
  \(U_{EM}F_{EM}=D\) recovers the monad.
\end{itemize}

The \textbf{counit of \(F_{EM}\dashv U_{EM}\) is exactly the algebra
structure map} --- \(\varepsilon_{(X,\alpha)}:(DX,\mu)\to(X,\alpha)\) is
\(\alpha\) itself. So calling the expectation \(\mathbb E\)
``counit-like'' in the previous group of turns was accurate:
\textbf{value mixing = a component of the counit of the EM adjunction}.

And the \textbf{comparison functor}
\(K:\mathrm{Kl}(D)\to\mathrm{EM}(D)\), \(X\mapsto(DX,\mu_X)\), is fully
faithful, embedding \(\mathrm{Kl}(D)\) as the \textbf{full subcategory
of free algebras}:

\[\mathrm{Kl}(D)\ \hookrightarrow\ \mathrm{EM}(D)\qquad(\text{free algebras}).\]

So the two are the ``free end (\(\mathrm{Kl}\), kernels)'' and the
``full-algebra end (\(\mathrm{EM}\), convex spaces)'': the value space
\(V\) lives in \(\mathrm{EM}(D)\) as a \textbf{non-free} \(D\)-algebra,
and the attention computation feeds the kernel \(A\) of
\(\mathrm{Kl}(D)\) into the structure map \(\alpha_V\) of \(V\), an
object of \(\mathrm{EM}(D)\) --- a move that \textbf{straddles the two
categories}. This is the \(D\)-side content of ``a Transformer is not a
single SMCC but a composite of a Markov category and an SMCC.''

\hypertarget{summary-table-and-implications-for-unification}{%
\subsubsection{Summary Table, and Implications for Unification}\label{summary-table-and-implications-for-unification}}

\begin{longtable}[]{@{}lll@{}}
\toprule
\begin{minipage}[b]{0.30\columnwidth}\raggedright
\strut
\end{minipage} & \begin{minipage}[b]{0.30\columnwidth}\raggedright
\(\mathrm{Kl}(D)\)\strut
\end{minipage} & \begin{minipage}[b]{0.30\columnwidth}\raggedright
\(\mathrm{EM}(D)=\mathrm{Alg}(D)\)\strut
\end{minipage}\tabularnewline
\midrule
\endhead
\begin{minipage}[t]{0.30\columnwidth}\raggedright
protagonist (morphism/object)\strut
\end{minipage} & \begin{minipage}[t]{0.30\columnwidth}\raggedright
stochastic kernel (morphism) \(X\to D(Y)\)\strut
\end{minipage} & \begin{minipage}[t]{0.30\columnwidth}\raggedright
\(D\)-algebra (object) \((X,\alpha)\)\strut
\end{minipage}\tabularnewline
\begin{minipage}[t]{0.30\columnwidth}\raggedright
concretely\strut
\end{minipage} & \begin{minipage}[t]{0.30\columnwidth}\raggedright
stochastic matrix (FinStoch)\strut
\end{minipage} & \begin{minipage}[t]{0.30\columnwidth}\raggedright
convex space / vector space (with expectation)\strut
\end{minipage}\tabularnewline
\begin{minipage}[t]{0.30\columnwidth}\raggedright
Transformer\strut
\end{minipage} & \begin{minipage}[t]{0.30\columnwidth}\raggedright
\textbf{softmax} = kernel \(A\)\strut
\end{minipage} & \begin{minipage}[t]{0.30\columnwidth}\raggedright
\textbf{value space} \(V\), \(\alpha_V=\mathbb E\) = value mixing\strut
\end{minipage}\tabularnewline
\begin{minipage}[t]{0.30\columnwidth}\raggedright
relation to the monad\strut
\end{minipage} & \begin{minipage}[t]{0.30\columnwidth}\raggedright
subcategory of free algebras\strut
\end{minipage} & \begin{minipage}[t]{0.30\columnwidth}\raggedright
full algebras, \(U_{EM}F_{EM}=D\)\strut
\end{minipage}\tabularnewline
\begin{minipage}[t]{0.30\columnwidth}\raggedright
copy\strut
\end{minipage} & \begin{minipage}[t]{0.30\columnwidth}\raggedright
non-natural (generates correlation)\strut
\end{minipage} & \begin{minipage}[t]{0.30\columnwidth}\raggedright
convex structure preserved by affine maps\strut
\end{minipage}\tabularnewline
\bottomrule
\end{longtable}

The unified category (composite monad \(M=T\circ D\))" also becomes
visible here. To bundle the linear part \(T\) and the distribution \(D\)
into a single monad, one needs a distributive law \(DT\Rightarrow TD\),
and its natural candidate is precisely \textbf{``expectation commutes
with linear maps''} --- the very fact that \(\alpha_V\) is
affine/linear. That the \(D\)-algebra structure is compatible with the
linear structure of \(V\) (\(\mathbb E\) is linear) is the seed that
generates the distributive law between \(D\) and \(T\). So pinning down
the relationship between \(\mathrm{Kl}(D)\) and \(\mathrm{EM}(D)\) leads
directly into the construction of the unified category (the contents of
the previous turn's Lean \texttt{composeMonad}).

Writing down the distributive law \(DT\Rightarrow TD\) from this
linearity of expectation = the \(D\)-algebra is the next move toward
unification and Lean formalization. If you want to proceed by concretely
defining the candidate for that natural transformation (corresponding to
``pushforward of expectation under a linear map,'' \(D(Tf)\to T(Df)\)),
you can write out each component of the distributive law by taking \(V\)
as the free vector space monad \(T\) and \(D\) as the distribution monad.

\end{document}